%!TEX root = main.tex

%\section{FIRRTL Syntax}

%%!TEX root = main.tex

\begin{figure}[htbp]
{\footnotesize
\[
\begin{array}{l l l}
%\mbox{\bf Circuit} & 
circuit & = & \mbox{``circuit''}, id, newline, indent, \{module \mid extmodule\}, dedent;\\
%
%\mbox{\bf Module} & 
module & = & \mbox{``module''} \ id, newline, indent, \\
%& 
& & \{port, newline\}, \{stmt, newline\}, dedent;\\
%
%\mbox{\bf External Module} & 
extmodule & = & \mbox{``extmodule''}, id, newline, indent,\\
%& 
  &  &  \{port, newline\}, [\mbox{``defname''}, \mbox{``=''}, id, newline], \\
% & 
 & & \{\mbox{``parameter''}, \mbox{``=''}, (string \mid int), newline\}, dedent;\\
%
%\mbox{\bf Port} & 
port & = & (\mbox{``input''} \mid \mbox{``output''}), id, \mbox{``:''}, type;\\
%
 %\mbox{\bf Statment} & 
 stmt & = & \mbox{``wire''}, id, \mbox{``:''}, type \\
 %& 
 & &  \mid \mbox{``reg''}, id, \mbox{``:''}, type, \mbox{``,''}, exp, [\mbox{``(with:\{reset=>(''}, exp, \mbox{``,''}, exp, \mbox{``)\})''}] \\
 %& 
  &   &  \mid memory \mid \mbox{``inst''}, id, \mbox{``of''}, id \mid \mbox{``node''}, id, \mbox{``=''}, id  \\
 %& 
  &   & \mid reference, \mbox{``<=''}, exp  \mid reference, \mbox{``<-''}, exp  \\
 %& 
  &   &\mid reference, \mbox{``is invalid''}  \mid \mbox{``attach(''}, \{reference\}, \mbox{``)''}  \\ 
%&  
&   &\mid \mbox{``when''}, exp, \mbox{``:''}, newline, indent, \{stmt\}, dedent \\
%& 
& & [\mbox{``else:''}, indent, \{stmt\}, dedent] \\ 
%&  
&   &\mid \mbox{``stop(''}, exp, \mbox{``,''}, exp, \mbox{``,''}, int, \mbox{``)''}\\ 
%&  
&   &\mid \mbox{``printf(''}, exp, \mbox{``,''}, exp, \mbox{``,''}, string, \mbox{``,''}, \{exp\}, \mbox{``)''}, [\mbox{``:''}, id]\\ 
%&  
&   &\mid \mbox{``skip''};\\ 
%
%\mbox{\bf Memory} & 
memory & = &  \mbox{``mem''}, id, newline, indent, \\
 %& 
 & & \mbox{``data-type''}, \mbox{``=>''}, type, newline, \\
 %& 
 & & \mbox{``depth''}, \mbox{``=>''}, int, newline, \\
 %& 
 & & \mbox{``read-latency''}, \mbox{``=>''}, int, newline, \\ 
 %& 
 & & \mbox{``write-latency}, \mbox{``=>''}, int, newline,  \\
 %& 
 & & \mbox{``read-under-write''}, \mbox{``=>''}, ruw, newline, \\
  %& 
  & & \{\mbox{``reader''}, \mbox{``=>''}, id, newline\}, \\
  %& 
  & & \{\mbox{``writer''}, \mbox{``=>''}, id, newline\}, \\
  %& 
  & & \{\mbox{``readwriter''}, \mbox{``=>''}, id, newline\}, dedent;\\    
%
%\mbox{\bf Read Under Write Flag} & 
ruw & = & \mbox{``old''} \mid \mbox{``new''} \mid \mbox{``undefined''};\\
%
%\mbox{\bf Expression} & 
exp & = & (\mbox{``UInt''} \mid \mbox{``SInt''}), [width], \mbox{``(''}, int, \mbox{``)''} \mid reference \\
%& 
& & \mid  \mbox{``mux''}, \mbox{``(''}, exp, \mbox{``,''} exp, \mbox{``,''}, exp, \mbox{``)''} \mid \mbox{``validif''}, \mbox{``(''}, exp, \mbox{``,''}, exp, \mbox{``)''}\\
%& 
& & \mid primop;\\
%
%& 
reference & = & id \mid reference, \mbox{``.''}, id \mid reference, \mbox{``[''}, int, \mbox{``]''} \mid reference, \mbox{``[''}, exp, \mbox{``]''};\\
%
%& 
width & = & \mbox{``<''}, int, \mbox{``>''};\\
%
%& 
binarypoint & = & \mbox{``<''}, \mbox{``<''}, int, \mbox{``>''}, \mbox{``>''};\\
%
 %\mbox{\bf Bundle Field} & 
 field & = & [\mbox{``flip''}]\ id: type;\\
%
%\mbox{\bf Data Types} & 
type & = & type\_ground \mid type\_aggregate; \\
 %& 
 type\_ground & = &  \mbox{``Clock''} \mid \mbox{``Reset''} \mid \mbox{``AsyncReset''} \mid (\mbox{``UInt''} \mid \mbox{``SInt''}  \mid \mbox{``Analog''}), [width]  \\
 %& 
 & & \mid \mbox{``Fixed''}, [width], [binarypoint];\\
 %& 
  type\_aggregate & = & \mbox{``\{''}, field, \{field\}, \mbox{``\}''} \mid type, \mbox{``[''}, int, \mbox{``]''};\\
%& & & \mbox{``UInt''}, [\mbox{``<''}, int, \mbox{``>''}] \mid \mbox{``SInt''}, [\mbox{``<''}, int, \mbox{``>''} ] \\
%& & & \mid \mbox{``Fixed''}, [\mbox{``<''}, int, \mbox{``>''} ] [\mbox{``<}\mbox{<''}, int, \mbox{``>}\mbox{>''}]\\
% &  &  & \mid \mbox{``Clock''} \mid \mbox{``Analog''},  [\mbox{``<''}, int, \mbox{``>''}] \mid \mbox{``\{''}, field, \mbox{``\}''} \mid type, \mbox{``[''}, int, \mbox{``]''}; \\
%
%\mbox{\bf Primitive Operation} & 
primop & = & primop\_2exp \mid primop\_1exp \mid primop\_1exp1int \mid primop\_1exp2int; \\
%& 
primop\_2exp & = & primop\_2exp\_keyword, \mbox{``(''} , exp, \mbox{``,''} , exp, \mbox{``)''}; \\
%
%& 
primop\_2exp\_keyword & = &  \mbox{``add''}  \mid  \mbox{``sub''}  \mid  \mbox{``sub''} \mid \mbox{``mul''} \mid \mbox{``div''} \mid \mbox{``mod''} \mid \mbox{``lt''} \mid \mbox{``leq''}  \\
%& 
& &  \mid \mbox{``gt''} \mid \mbox{``geq''} \mid \mbox{``eq''} \mid \mbox{``neq''} \mid \mbox{``dshl''}  \mid \mbox{``dshr''} \\
%& 
 & & \mid \mbox{``and''} \mid \mbox{``or''} \mid \mbox{``xor''} \mid \mbox{``cat''}; \\
 %
%& 
primop\_1exp & = & primop\_1exp\_keyword, \mbox{``(''}, exp, \mbox{``)''}; \\
%& 
primop\_1exp\_keyword & = & \mbox{``asUInt''} \mid \mbox{``asSInt''} \mid \mbox{``asClock''} \mid \mbox{``cvt''} \mid \mbox{``neg''} \mid \mbox{``not''}  \\
%& 
& & \mid \mbox{``andr''} \mid \mbox{``orr''} \mid  \mbox{``xorr''};\\
%
%& 
primop\_1exp1int & = & primop\_1exp1int\_keyword, \mbox{``(''}, exp, \mbox{``,''}, int, \mbox{``)''}; \\
%& 
primop\_1exp1int\_keyword & = & \mbox{``pad''} \mid \mbox{``shl''} \mid \mbox{``shr''} \mid \mbox{``head''} \mid \mbox{``tail''}; \\
%& 
primop\_1exp2int & = & primop\_1exp2int\_keyword, \mbox{``(''}, exp, \mbox{``,''}, int, \mbox{``,''}, int \mbox{``)''}; \\
%& 
primop\_1exp2int\_keyword & = & \mbox{``bits''}; \\
%
%\mbox{\bf File Information Token} & 
%info & = & \mbox{``@''}, \mbox{``[''}, \{string, \mbox{`` ''}, linecol\}, \mbox{``]''};\\
%& 
linecol & = &  digit\_dec, \{digit\_dec\}, \mbox{``:''}, digit\_dec, \{digit\_dec\};\\
%& 
id & = & ( \mbox{``\_''} \mid letter), \{ \mbox{``\_''} \mid letter \mid digit\_dec\};
%
\end{array}
\]
}
\caption{Syntax of FIRRTL}
\label{fig-firrtl-synt}
\end{figure}



%%%%%%%%%%%%%%%%%%%%%%%%%%%%%%%%%%%%%%%%%%%%%%%
%%%%%%%%%%%%%%%%%%syntax with info %%%%%%%%%%%%%%%%%%%
%%%%%%%%%%%%%%%%%%%%%%%%%%%%%%%%%%%%%%%%%%%%%%%

\hide{

\begin{figure}[t]
{\footnotesize
\[
\begin{array}{l l l}
%\mbox{\bf Circuit} & 
circuit & = & \mbox{``circuit''}, id, \mbox{``:''}, [info], newline, indent, \{module \mid extmodule\}, dedent;\\
%
%\mbox{\bf Module} & 
module & = & \mbox{``module''} \ id, \mbox{``:''}, [info], newline, indent, \\
%& 
& & \{port, newline\}, \{stmt, newline\}, dedent;\\
%
%\mbox{\bf External Module} & 
extmodule & = & \mbox{``extmodule''}, id, \mbox{``:''}, [info], newline, indent,\\
%& 
  &  &  \{port, newline\}, [\mbox{``defname''}, \mbox{``=''}, id, newline], \\
% & 
 & & \{\mbox{``parameter''}, \mbox{``=''}, (string \mid int), newline\}, dedent;\\
%
%\mbox{\bf Port} & 
port & = & (\mbox{``input''} \mid \mbox{``output''}), id, \mbox{``:''}, type, [info];\\
%
 %\mbox{\bf Statment} & 
 stmt & = & \mbox{``wire''}, id, \mbox{``:''}, type, [info] \\
 %& 
 & &  \mid \mbox{``reg''}, id, \mbox{``:''}, type, \mbox{``,''}, exp, [\mbox{``(with:\{reset=>(''}, exp, \mbox{``,''}, exp, \mbox{``)\})''}], [info] \\
 %& 
  &   &  \mid memory \mid \mbox{``inst''}, id, \mbox{``of''}, id, [info] \mid \mbox{``node''}, id, \mbox{``=''}, id, [info]  \\
 %& 
  &   & \mid reference, \mbox{``<=''}, exp, [info]  \mid reference, \mbox{``<-''}, exp, [info]  \\
 %& 
  &   &\mid reference, \mbox{``is invalid''}, [info]  \mid \mbox{``attach(''}, \{reference\}, \mbox{``)''}, [info]  \\ 
%&  
&   &\mid \mbox{``when''}, exp, \mbox{``:''}, [info], newline, indent, \{stmt\}, dedent \\
%& 
& & [\mbox{``else:''}, indent, \{stmt\}, dedent] \\ 
%&  
&   &\mid \mbox{``stop(''}, exp, \mbox{``,''}, exp, \mbox{``,''}, int, \mbox{``)''}, [info]\\ 
%&  
&   &\mid \mbox{``printf(''}, exp, \mbox{``,''}, exp, \mbox{``,''}, string, \mbox{``,''}, \{exp\}, \mbox{``)''}, [\mbox{``:''}, id], [info]\\ 
%&  
&   &\mid \mbox{``skip''}, [info];\\ 
%
%\mbox{\bf Memory} & 
memory & = &  \mid  \mbox{``mem''}, id, \mbox{``:''}, [info], newline, indent, \\
 %& 
 & & \mbox{``data-type''}, \mbox{``=>''}, type, newline, \\
 %& 
 & & \mbox{``depth''}, \mbox{``=>''}, int, newline, \\
 %& 
 & & \mbox{``read-latency''}, \mbox{``=>''}, int, newline, \\ 
 %& 
 & & \mbox{``write-latency}, \mbox{``=>''}, int, newline,  \\
 %& 
 & & \mbox{``read-under-write''}, \mbox{``=>''}, ruw, newline, \\
  %& 
  & & \{\mbox{``reader''}, \mbox{``=>''}, id, newline\}, \\
  %& 
  & & \{\mbox{``writer''}, \mbox{``=>''}, id, newline\}, \\
  %& 
  & & \{\mbox{``readwriter''}, \mbox{``=>''}, id, newline\}, dedent;\\    
%
%\mbox{\bf Read Under Write Flag} & 
ruw & = & \mbox{``old''} \mid \mbox{``new''} \mid \mbox{``undefined''};\\
%
%\mbox{\bf Expression} & 
exp & = & (\mbox{``UInt''} \mid \mbox{``SInt''}), [width], \mbox{``(''}, int, \mbox{``)''} \mid reference \\
%& 
& & \mid  \mbox{``mux''}, \mbox{``(''}, exp, \mbox{``,''} exp, \mbox{``,''}, exp, \mbox{``)''} \mid \mbox{``validif''}, \mbox{``(''}, exp, \mbox{``,''}, exp, \mbox{``)''}\\
%& 
& & \mid primop;\\
%
%& 
reference & = & id \mid reference, \mbox{``.''}, id \mid reference, \mbox{``[''}, int, \mbox{``]''} \mid reference, \mbox{``[''}, exp, \mbox{``]''};\\
%
%& 
width & = & \mbox{``<''}, int, \mbox{``>''};\\
%
%& 
binarypoint & = & \mbox{``<''}, \mbox{``<''}, int, \mbox{``>''}, \mbox{``>''};\\
%
 %\mbox{\bf Bundle Field} & 
 field & = & [\mbox{``flip''}]\ id: type;\\
%
%\mbox{\bf Data Types} & 
type & = & type\_ground \mid type\_aggregate; \\
 %& 
 type\_ground & = &  \mbox{``Clock''} \mid \mbox{``Reset''} \mid \mbox{``AsyncReset''} \mid (\mbox{``UInt''} \mid \mbox{``SInt''}  \mid \mbox{``Analog''}), [width]  \\
 %& 
 & & \mid \mbox{``Fixed''}, [width], [binarypoint];\\
 %& 
  type\_aggregate & = & \mbox{``\{''}, field, \{field\}, \mbox{``\}''} \mid type, \mbox{``[''}, int, \mbox{``]''};\\
%& & & \mbox{``UInt''}, [\mbox{``<''}, int, \mbox{``>''}] \mid \mbox{``SInt''}, [\mbox{``<''}, int, \mbox{``>''} ] \\
%& & & \mid \mbox{``Fixed''}, [\mbox{``<''}, int, \mbox{``>''} ] [\mbox{``<}\mbox{<''}, int, \mbox{``>}\mbox{>''}]\\
% &  &  & \mid \mbox{``Clock''} \mid \mbox{``Analog''},  [\mbox{``<''}, int, \mbox{``>''}] \mid \mbox{``\{''}, field, \mbox{``\}''} \mid type, \mbox{``[''}, int, \mbox{``]''}; \\
%
%\mbox{\bf Primitive Operation} & 
primop & = & primop\_2exp \mid primop\_1exp \mid primop\_1exp1int \mid primop\_1exp2int; \\
%& 
primop\_2exp & = & primop\_2exp\_keyword, \mbox{``(''} , exp, \mbox{``,''} , exp, \mbox{``)''}; \\
%
%& 
primop\_2exp\_keyword & = &  \mbox{``add''}  \mid  \mbox{``sub''}  \mid  \mbox{``sub''} \mid \mbox{``mul''} \mid \mbox{``div''} \mid \mbox{``mod''} \mid \mbox{``lt''} \mid \mbox{``leq''}  \\
%& 
& &  \mid \mbox{``gt''} \mid \mbox{``geq''} \mid \mbox{``eq''} \mid \mbox{``neq''} \mid \mbox{``dshl''}  \mid \mbox{``dshr''} \\
%& 
 & & \mid \mbox{``and''} \mid \mbox{``or''} \mid \mbox{``xor''} \mid \mbox{``cat''}; \\
 %
%& 
primop\_1exp & = & primop\_1exp\_keyword, \mbox{``(''}, exp, \mbox{``)''}; \\
%& 
primop\_1exp\_keyword & = & \mbox{``asUInt''} \mid \mbox{``asSInt''} \mid \mbox{``asClock''} \mid \mbox{``cvt''} \mid \mbox{``neg''} \mid \mbox{``not''}  \\
%& 
& & \mid \mbox{``andr''} \mid \mbox{``orr''} \mid  \mbox{``xorr''};\\
%
%& 
primop\_1exp1int & = & primop\_1exp1int\_keyword, \mbox{``(''}, exp, \mbox{``,''}, int, \mbox{``)''}; \\
%& 
primop\_1exp1int\_keyword & = & \mbox{``pad''} \mid \mbox{``shl''} \mid \mbox{``shr''} \mid \mbox{``head''} \mid \mbox{``tail''}; \\
%& 
primop\_1exp2int & = & primop\_1exp2int\_keyword, \mbox{``(''}, exp, \mbox{``,''}, int, \mbox{``,''}, int \mbox{``)''}; \\
%& 
primop\_1exp2int\_keyword & = & \mbox{``bits''}; \\
%
%\mbox{\bf File Information Token} & 
info & = & \mbox{``@''}, \mbox{``[''}, \{string, \mbox{`` ''}, linecol\}, \mbox{``]''};\\
%& 
linecol & = &  digit\_dec, \{digit\_dec\}, \mbox{``:''}, digit\_dec, \{digit\_dec\};\\
%& 
id & = & ( \mbox{``\_''} \mid letter), \{ \mbox{``\_''} \mid letter \mid digit\_dec\};
%
\end{array}
\]
}
\caption{Syntax of FIRRTL}
\label{fig-firrtl-synt}
\end{figure}


}

\section{Missing details in Section~\ref{sec-form-firrtl}}\label{app-lofirrtl}


\zhilin{define $\typeofexp(e, tmap)$ here}

\paragraph*{The definitions of $te(e)$ and $s(e)$}
%
For $\theta \bwidth{n}(n')$ such that $\theta \in \{\mbox{\prettyll{UInt}, \prettyll{SInt}}\}$, $n > 0$, and $n'$ is an integer, we define $bv_{\theta \bwidth{n}(n')}$ as follows.
\begin{itemize}
\item For $n > 0$ and $n' \ge 0$ such that $n' < 2^n$, let $bv_{\mbox{\prettyll{UInt}}<n>(n')}$ denote the bit vector of length $n$ representing the unsigned integer $n'$.
%
\item For $n > 0$ and an integer $n'$ such that $-2^{n-1} \le n' < 2^{n-1}$, let $bv_{\mbox{\prettyll{SInt}}<n>(n')}$ denote the bit vector of length $n$ representing the signed integer $n'$.
\end{itemize}
The function $te$ and $s$ from $\cpnts(A)$ to $\gtypes$ are extended to expressions as follows.
\begin{itemize}
\item $te(\theta \bwidth{n}(c)) = \theta \bwidth{n}$ and $s(\theta \bwidth{n}(c)) = bv_{\theta \bwidth{n}(c)}$, where $\theta \in \{\mbox{\prettyll{UInt}}, \mbox{\prettyll{SInt}}\}$,
%\item $te(\sint{n}(c)) = \sint{n}$,
%
\item let $te(e_1) = \theta \bwidth{n_1}$ and $te(e_2) = \theta \bwidth{n_2}$, where $\theta \in \{\mbox{\prettyll{UInt}}, \mbox{\prettyll{SInt}}\}$, then $te(mux(e_0, e_1, e_2)) = \theta\bwidth{\max(n_1, n_2)}$, moreover, if $s(e_0) = 1$, then $s(mux(e_0, e_1, e_2)) = bv_{\theta\bwidth{n_0}(s(e_1))}$, otherwise, $s(mux(e_0, e_1, e_2)) = bv_{\theta\bwidth{n_0}(s(e_2))}$,
%
\item $te(validif(e_1, e_2)) = te(e_2)$, moreover, if $s(e_1) = 1$, then $s(validif(e_1, e_2)) = s(e_2)$, otherwise, $s(validif(e_1, e_2))  = \bot$,
%
\item let $te(e_1) = \theta \bwidth{n_1}$ and $te(e_2) = \theta \bwidth{n_2}$, where $\theta \in \{\mbox{\prettyll{UInt}}, \mbox{\prettyll{SInt}}\}$, then
\begin{itemize}
\item $te(add(e_1, e_2)) =\theta \bwidth{\max(n_1, n_2)+1}$ and $te(sub(e_1, e_2)) = \theta \bwidth{\max(n_1, n_2)+1}$, moreover, $s(add(e_1, e_2)) = bv_{\theta \bwidth{\max(n_1,n_2)+1} (s(e_1) + s(e_2))}$ and \\
$s(sub(e_1, e_2)) = bv_{\theta \bwidth{\max(n_1,n_2)+1} (s(e_1) - s(e_2))}$,
%
\item $te(mul(e_1, e_2)) = \theta \bwidth{n_1+n_2}$ and $s(mul(e_1, e_2)) = bv_{\theta \bwidth{n_1+n_2} (s(e_1) \times s(e_2))}$,
%
\item  if $\theta = \mbox{\prettyll{UInt}}$, then $te(div(e_1, e_2)) =\mbox{\prettyll{UInt}} \bwidth{n_1}$ and $s(div(e_1, e_2)) = bv_{\mbox{\prettyll{UInt}} \bwidth{n_1} (s(e_1) \div s(e_2))}$, otherwise, $te(div(e_1, e_2)) = \mbox{\prettyll{SInt}} \bwidth{n_1+1}$ and $s(div(e_1, e_2)) = bv_{\mbox{\prettyll{SInt}} \bwidth{n_1+1} (s(e_1) \div s(e_2))}$,
%
\item $te(mod(e_1, e_2)) = \theta \bwidth{\min(n_1,n_2)}$ and $s(mod(e_1, e_2)) = bv_{\theta \bwidth{\min(n_1,n_2)} (s(e_1) \bmod s(e_2))}$,
%
\item $te(cp(e_1, e_2)) = \mbox{\prettyll{UInt}} \bwidth{1}$, moreover, $s(cp(e_1, e_2)) = 1$ iff $s(e_1)\ cp\ s(e_2)$, \\ where $cp \in \{lt, leq, gt, geq, eq, neq\}$,
%
\item $te(dshl(e_1, e_2)) = \theta \bwidth{n_1+2^{n_2}-1}$ and $s(dshl(e_1, e_2)) = bv_{\theta \bwidth{n_1+2^{n_2}-1}(s(e_1) \times 2^{s(e_2)})}$,
%
\item $te(dshr(e_1, e_2)) = \theta \bwidth{n_1}$ and $s(dshr(e_1, e_2)) = bv_{\theta \bwidth{n_1}(s(e_1) \div 2^{s(e_2)})}$,
%
\item $te(and(e_1, e_2)) = te(or(e_1, e_2)) =  te(xor(e_1, e_2)) = \mbox{\prettyll{UInt}} \bwidth{\max(n_1, n_2)}$, moreover,
$s(and(e_1, e_2))$ is the bitwise and of $bv_{\theta\bwidth{\max(n_1, n_2)}(s(e_1))}$ and $bv_{\theta\bwidth{\max(n_1, n_2)}(s(e_2))}$,  $s(or(e_1, e_2))$ is the bitwise or of $bv_{\theta\bwidth{\max(n_1, n_2)}(s(e_1))}$ and $bv_{\theta\bwidth{\max(n_1, n_2)}(s(e_2))}$, $s(xor(e_1, e_2))$ is the bitwise xor of $bv_{\theta\bwidth{\max(n_1, n_2)}(s(e_1))}$ and $bv_{\theta\bwidth{\max(n_1, n_2)}(s(e_2))}$,
%
\item $te(cat(e_1, e_2)) = \mbox{\prettyll{UInt}} \bwidth{n_1 + n_2}$, moreover,
\begin{itemize}
\item if either $\theta = \mbox{\prettyll{UInt}}$, or $\theta=\mbox{\prettyll{SInt}}$, $s(e_1) \ge 0$, and $s(e_2) \ge 0$, then \\ $s(cat(e_1, e_2))$ is $bv_{\mbox{\prettyll{UInt}}\bwidth{n_1+n_2}(s(e_1) \times 2^{n_2} + s(e_2))}$,
%
\item if  $\theta=\mbox{\prettyll{SInt}}$, $s(e_1) \ge 0$, and $s(e_2) < 0$, then \\ $s(cat(e_1, e_2))$ is $bv_{\mbox{\prettyll{UInt}}\bwidth{n_1+n_2}(s(e_1) \times 2^{n_2} + 2^{n_2} + s(e_2))}$,
%
\item if  $\theta=\mbox{\prettyll{SInt}}$ and $s(e_1) < 0, s(e_2) \ge 0$, then \\  $s(cat(e_1, e_2))$ is $bv_{\mbox{\prettyll{UInt}}\bwidth{n_1+n_2}((2^{n_1}+s(e_1)) \times 2^{n_2} + s(e_2))}$,
%
\item if  $\theta=\mbox{\prettyll{SInt}}$ and $s(e_1) < 0, s(e_2) < 0$ then \\  $s(cat(e_1, e_2))$ is $bv_{\mbox{\prettyll{UInt}}\bwidth{n_1+n_2}((2^{n_1}-s(e_1)) \times 2^{n_2} + 2^{n_2}-s(e_2))}$,
\end{itemize}
\end{itemize}
%
\item let $te(e_1) = \theta \bwidth{n_1}$ and $te(e_2) = \mbox{\prettyll{UInt}} \bwidth{n_2}$, where $\theta \in \{\mbox{\prettyll{UInt}}, \mbox{\prettyll{SInt}}\}$, then
\begin{itemize}
\item $te(dshl(e_1, e_2)) = \theta \bwidth{n_1+2^{n_2}-1}$ and $s(dshl(e_1, e_2)) = bv_{\theta \bwidth{n_1+2^{n_2}-1}(s(e_1) \times 2^{s(e_2)})}$,
%
\item $te(dshr(e_1, e_2)) = \theta \bwidth{n_1}$ and $s(dshr(e_1, e_2)) = bv_{\theta \bwidth{n_1}(s(e_1) \div 2^{s(e_2)})}$,
%
\end{itemize}
%
\item let $te(e_1) = \theta \bwidth{n_1}$, where $\theta \in \{\mbox{\prettyll{UInt}}, \mbox{\prettyll{SInt}}\}$, then
\begin{itemize}
\item $te(andr(e_1)) = te(orr(e_1)) =  te(xorr(e_1)) = \uint{1}$, moreover, $s(andr(e_1))$ (resp. $s(orr(e_1))$, $s(xorr(e_1))$) is the and (resp. or, xor) of all the bits in the bit vector $s(e_1)$,
%
\item $te(neg(e_1)) = \sint{n_1 +1}$, $te(not(e_1)) =  \uint{n_1}$, moreover,
\begin{itemize}
\item $s(neg(e_1)) = bv_{\mbox{\prettyll{SInt}}\bwidth{n_1+1}(-s(e_1))}$,
%
\item $s(not(e_1))$ is obtained by applying the logical not to each bit of the bit vector $s(e_1)$,
\end{itemize}
%\begin{itemize}
%\item if $\theta = UInt$, then $s(neg(e_1)) = bv_{SInt\bwidth{n_1+1}(-s(e_1))}$,
%
%\item if $\theta = SInt$, then $s(neg(e_1)) = bv_{SInt\bwidth{n_1+1}(-s(e_1))}$.
%\end{itemize}
%
\item $te(cvt(e_1)) = \sint{n_1 +1}$ if $\theta = \mbox{\prettyll{UInt}}$, and $te(cvt(e_1)) =  \sint{n_1}$ otherwise,  moreover,
\begin{itemize}
\item if $\theta = \mbox{\prettyll{UInt}}$, then $s(cvt(e_1)) = bv_{\mbox{\prettyll{SInt}}\bwidth{n_1+1}(s(e_1))}$,
%
\item if $\theta = \mbox{\prettyll{SInt}}$, then $s(cvt(e_1)) = s(e_1)$,
\end{itemize}
\end{itemize}
%
\item let $te(e_1) = \theta \bwidth{n_1}$, where $\theta \in \{\mbox{\prettyll{UInt}}, \mbox{\prettyll{SInt}}\}$, then
\begin{itemize}
\item $te(pad(e_1, n)) = \theta\bwidth{\max(n_1, n)}$, moreover, if $n_1 < n$, then $s(pad(e_1,n)) = bv_{\theta\bwidth{n}(s(e_1))}$, otherwise, $s(pad(e_1,n)) = s(e_1)$,
%
\item $te(shl(e_1,n)) = \theta \bwidth{n_1 + n}$, moreover, $s(shl(e_1,n)) =  bv_{\theta\bwidth{n_1 + n}(2^{n} \times s(e_1))}$,
%\begin{itemize}
%\item if $\theta = UInt$, then $s(shl(e_1,n)) =  bv_{UInt\bwidth{n_1 + n}(2^{n} \times s(e_1))}$,
%\item if $\theta = SInt$, then $s(shl(e_1,n)) = bv_{SInt\bwidth{n_1 + n}(2^{n} \times s(e_1))}$,
%\end{itemize}
%
\item $te(shr(e_1, n)) = \theta \bwidth{n_1-n}$ if $n_1 > n$, and $te(shr(e_1, n)) =  \theta \bwidth{1}$ otherwise,  moreover,
\begin{itemize}
\item if $n_1 > n$, then $s(shr(e_1, n))$ is obtained from $s(e_1)$ by shifting $n$ bits right,
%
\item otherwise, if $\theta = \mbox{\prettyll{UInt}}$, then $s(shr(e_1, n))=bv_{\mbox{\prettyll{UInt}}<1>(0)}$, otherwise, $s(shr(e_1, n))=bv_{\mbox{\prettyll{SInt}}<1>(0)}$ if $s(e_1) \ge 0$, and $s(shr(e_1, n))=bv_{\mbox{\prettyll{SInt}}<1>(1)}$ if $s(e_1) < 0$,
\end{itemize}
%
\item $te(head(e_1, n)) = { \mbox{\prettyll{UInt}}}{n}$, $te(tail(e_1, n)) = { \mbox{\prettyll{UInt}}}{n_1-n}$, moreover, $s(head(e_1, n))$ is the most significant $n$ bits of $s(e_1)$, $s(tail(e_1, n))$ is obtained from $s(e_1)$ by truncating the $n$ most significant bits of $s(e_1)$,
%
\item $te(bits(e_1, n, n')) = \uint{n'-n+1}$, and $s(bits(e_1, n, n'))$ is the bit vector comprising the $n, n+1, \cdots, n'$-th bit of $s(e_1)$.
\end{itemize}
%
\end{itemize}


\paragraph*{The definition of ${\tt type\_of\_expr}(e)$ }%\zhilin{@Xiaomu, please add the definition here}
We define the type of HiFIRRTL program in the Coq as $ftype$. In general, it can be used to represent either a {\it ground type} ($gtyp$), a {\it vector type} ($vtyp$), or a {\tt bundle type} ($btyp$),
$$ftype::=gtyp\;|\;vtyp\,(ftype\,n)\;|\;btyp\;fields$$
where $$fields::=fnil\;|\;field\;fid\;flip\;ftype\;fields.$$
For each kind of expression of HiFIRRRL, we derived its type $t$ in $ftype$ as follow.
\begin{itemize}
\item ${\tt type\_of\_expr}(\mbox{\prettyll{id}})=t$, where reference $id$ is the entry in the component environment $ce$.
\item ${\tt type\_of\_expr}(\mbox{\prettyll{id.fid}})$ requires to iteratively match the type of $id$ to find the same $fid$ and then to return the $ftype$ associated. If the field identifier $fid$ can not be found in the $btyp$ of the reference $id$ or the type of $id$ is not $btyp$, it returns $\bot$.
\item ${\tt type\_of\_expr}(\mbox{\prettyll{id[n]}})$ returns the base $ftype$ of reference $id$. If the nature number $n=0$ or the type of reference $id$ is other than a $vtyp$, it returns $\bot$.
\item let ${\tt type\_of\_expr}(e_1)=t_1$ and ${\tt type\_of\_expr}(e_2)=t_2$, then ${\tt type\_of\_expr}(\mbox{\prettyll{mux}}(e_0,e_1,e_2))={\tt max\_width}(e_1,e_2)$, where the {\tt max\_width} compute recursively the max width of the (sub) components, iff $t_1$ $t_2$ are type equivalent, otherwise it returns $\bot$.
\item ${\tt type\_of\_expr}(\mbox{\prettyll{validif}}(e_1,e_2))={\tt type\_of\_expr}(e_2)$, regardless the type of $e1$.
\item ${\tt type\_of\_expr}(op)$, $op$ includes all the binary and unary operations, in which only $gtyp$ operation will involve. The {\tt type\_of\_expr} computes the type as same as $te$.
\end{itemize}



%!TEX root = main.tex

%%%%%%%%%%%%%%%%%%%%%%%%%%%%%%%%%%%%%%%%%%%%%%%%%%%%%%%%%%%%%%%%%%%%%%%%%%%%%% 
%\tl{stuff from here can be moved to appendix, they are just tedious defs.}
%%%%%%%%%%%%%%%%%%%%%%%%%%%%%%%%%%%%%%%%%%%%%%%%%%%%%%%%%%%%%%%%%%%%%%%%%%%%%%%
%%%%% vm_match

In this section, we present the missing definitions of functions and relations in Section~\ref{sec-form-firrtl}.

Let us first define $\typeofexp(e, tmap)$.
%
%!TEX root = main.tex

\begin{itemize}
\item $\typeofexp(x, tmap)$ for name $x$ in $M$ is defined as follows. 
\begin{itemize}
  \item If $tmap(id(x)) = (\uintimp, n)$ or $(\sintimp, n)$ for some $n \in \Nat$, then the $\typeofexp(x, tmap)$ is defined as $(\uintimp, n')$ or $(\sintimp, n')$ where $n' \in \Nat$ is nondeterministically chosen. This nondeterministic choice is safe since expressions are used on the right-hand side of connect statements and all the chosen widths have to satisfy the {\typecompatible} constraints. 
%
\item If $tmap(id(x)) \in \gtypes \setminus \{(\uintimp, n), (\sintimp, n) \mid n \in \Nat\}$, then $\typeofexp(x, tmap) = tmap(id(x))$.
%
\item If $tmap(id(x))$ is an aggregate type, then {\typeofexp} inductively make nondeterministic choices for all the ground types of implicit widths in it. 
\end{itemize}
%
\item $\typeofexp(x.f, tmap)$ 
  is inductively defined as follows: Suppose $\typeofexp(x, tmap) = (\btype, ft)$. If $f$ occurs in $ft$, that is, $ft$ contains a triple $(f, b, t)$, then $\typeofexp(x.f, tmap)$ is defined as $t$, otherwise, $\typeofexp(x.f, tmap) = \svnone$. 
%
\item $\typeofexp(r[n], tmap)$ is inductively defined as follows: Suppose $\typeofexp(r, tmap) = (\vtype, t[m])$ for some $m \in \Nat$. If $n < m$, then $\typeofexp(r[n], tmap) = t$, otherwise, $\typeofexp(r[n], tmap) = \svnone$. 
%
\item $\typeofexp(\mbox{\prettyll{mux}}(c, e_1,e_2), tmap)= t$ is defined as follows: If $\typeofexp(c, tmap)= (\uint, 1)$ or $(\uintimp, 1)$, and $t_1 = \typeofexp(e_1, tmap)$ is equivalent to $t_2 = \typeofexp(e_2, tmap)$ after ignoring the widths (e.g. $\uint$ is equivalent to $\uint$, but not $\sint$), then $t$ is obtained from $t_1, t_2$ by taking the maximum of the widths. 
%
\item $\typeofexp(uop\ e_1)$ and $\typeofexp(e_1 \ bop\ e_2)$ for unary operations $uop$ and binary operations $bop$ can be defined according to the semantics of these operations. We  omit their definitions here, to avoid tediousness.
\end{itemize}


%%%%%%%%%%%%%%
%%%%%%%%%%%%%%
\hide{
%
For $\theta \bwidth{n}(n')$ such that $\theta \in \{\mbox{\prettyll{UInt}, \prettyll{SInt}}\}$, $n > 0$, and $n'$ is an integer, we define $bv_{\theta \bwidth{n}(n')}$ as follows.
\begin{itemize}
\item For $n > 0$ and $n' \ge 0$ such that $n' < 2^n$, let $bv_{\mbox{\prettyll{UInt}}<n>(n')}$ denote the bit vector of length $n$ representing the unsigned integer $n'$.
%
\item For $n > 0$ and an integer $n'$ such that $-2^{n-1} \le n' < 2^{n-1}$, let $bv_{\mbox{\prettyll{SInt}}<n>(n')}$ denote the bit vector of length $n$ representing the signed integer $n'$.
\end{itemize}
The function $te$ and $s$ from $\cpnts(A)$ to $\gtypes$ are extended to expressions as follows.
\begin{itemize}
\item $te(\theta \bwidth{n}(c)) = \theta \bwidth{n}$ and $s(\theta \bwidth{n}(c)) = bv_{\theta \bwidth{n}(c)}$, where $\theta \in \{\mbox{\prettyll{UInt}}, \mbox{\prettyll{SInt}}\}$,
%\item $te(\sint{n}(c)) = \sint{n}$,
%
\item let $te(e_1) = \theta \bwidth{n_1}$ and $te(e_2) = \theta \bwidth{n_2}$, where $\theta \in \{\mbox{\prettyll{UInt}}, \mbox{\prettyll{SInt}}\}$, then $te(mux(e_0, e_1, e_2)) = \theta\bwidth{\max(n_1, n_2)}$, moreover, if $s(e_0) = 1$, then $s(mux(e_0, e_1, e_2)) = bv_{\theta\bwidth{n_0}(s(e_1))}$, otherwise, $s(mux(e_0, e_1, e_2)) = bv_{\theta\bwidth{n_0}(s(e_2))}$,
%
\item $te(validif(e_1, e_2)) = te(e_2)$, moreover, if $s(e_1) = 1$, then $s(validif(e_1, e_2)) = s(e_2)$, otherwise, $s(validif(e_1, e_2))  = \bot$,
%
\item let $te(e_1) = \theta \bwidth{n_1}$ and $te(e_2) = \theta \bwidth{n_2}$, where $\theta \in \{\mbox{\prettyll{UInt}}, \mbox{\prettyll{SInt}}\}$, then
\begin{itemize}
\item $te(add(e_1, e_2)) =\theta \bwidth{\max(n_1, n_2)+1}$ and $te(sub(e_1, e_2)) = \theta \bwidth{\max(n_1, n_2)+1}$, moreover, $s(add(e_1, e_2)) = bv_{\theta \bwidth{\max(n_1,n_2)+1} (s(e_1) + s(e_2))}$ and \\
$s(sub(e_1, e_2)) = bv_{\theta \bwidth{\max(n_1,n_2)+1} (s(e_1) - s(e_2))}$,
%
\item $te(mul(e_1, e_2)) = \theta \bwidth{n_1+n_2}$ and $s(mul(e_1, e_2)) = bv_{\theta \bwidth{n_1+n_2} (s(e_1) \times s(e_2))}$,
%
\item  if $\theta = \mbox{\prettyll{UInt}}$, then $te(div(e_1, e_2)) =\mbox{\prettyll{UInt}} \bwidth{n_1}$ and $s(div(e_1, e_2)) = bv_{\mbox{\prettyll{UInt}} \bwidth{n_1} (s(e_1) \div s(e_2))}$, otherwise, $te(div(e_1, e_2)) = \mbox{\prettyll{SInt}} \bwidth{n_1+1}$ and $s(div(e_1, e_2)) = bv_{\mbox{\prettyll{SInt}} \bwidth{n_1+1} (s(e_1) \div s(e_2))}$,
%
\item $te(mod(e_1, e_2)) = \theta \bwidth{\min(n_1,n_2)}$ and $s(mod(e_1, e_2)) = bv_{\theta \bwidth{\min(n_1,n_2)} (s(e_1) \bmod s(e_2))}$,
%
\item $te(cp(e_1, e_2)) = \mbox{\prettyll{UInt}} \bwidth{1}$, moreover, $s(cp(e_1, e_2)) = 1$ iff $s(e_1)\ cp\ s(e_2)$, \\ where $cp \in \{lt, leq, gt, geq, eq, neq\}$,
%
\item $te(dshl(e_1, e_2)) = \theta \bwidth{n_1+2^{n_2}-1}$ and $s(dshl(e_1, e_2)) = bv_{\theta \bwidth{n_1+2^{n_2}-1}(s(e_1) \times 2^{s(e_2)})}$,
%
\item $te(dshr(e_1, e_2)) = \theta \bwidth{n_1}$ and $s(dshr(e_1, e_2)) = bv_{\theta \bwidth{n_1}(s(e_1) \div 2^{s(e_2)})}$,
%
\item $te(and(e_1, e_2)) = te(or(e_1, e_2)) =  te(xor(e_1, e_2)) = \mbox{\prettyll{UInt}} \bwidth{\max(n_1, n_2)}$, moreover,
$s(and(e_1, e_2))$ is the bitwise and of $bv_{\theta\bwidth{\max(n_1, n_2)}(s(e_1))}$ and $bv_{\theta\bwidth{\max(n_1, n_2)}(s(e_2))}$,  $s(or(e_1, e_2))$ is the bitwise or of $bv_{\theta\bwidth{\max(n_1, n_2)}(s(e_1))}$ and $bv_{\theta\bwidth{\max(n_1, n_2)}(s(e_2))}$, $s(xor(e_1, e_2))$ is the bitwise xor of $bv_{\theta\bwidth{\max(n_1, n_2)}(s(e_1))}$ and $bv_{\theta\bwidth{\max(n_1, n_2)}(s(e_2))}$,
%
\item $te(cat(e_1, e_2)) = \mbox{\prettyll{UInt}} \bwidth{n_1 + n_2}$, moreover,
\begin{itemize}
\item if either $\theta = \mbox{\prettyll{UInt}}$, or $\theta=\mbox{\prettyll{SInt}}$, $s(e_1) \ge 0$, and $s(e_2) \ge 0$, then \\ $s(cat(e_1, e_2))$ is $bv_{\mbox{\prettyll{UInt}}\bwidth{n_1+n_2}(s(e_1) \times 2^{n_2} + s(e_2))}$,
%
\item if  $\theta=\mbox{\prettyll{SInt}}$, $s(e_1) \ge 0$, and $s(e_2) < 0$, then \\ $s(cat(e_1, e_2))$ is $bv_{\mbox{\prettyll{UInt}}\bwidth{n_1+n_2}(s(e_1) \times 2^{n_2} + 2^{n_2} + s(e_2))}$,
%
\item if  $\theta=\mbox{\prettyll{SInt}}$ and $s(e_1) < 0, s(e_2) \ge 0$, then \\  $s(cat(e_1, e_2))$ is $bv_{\mbox{\prettyll{UInt}}\bwidth{n_1+n_2}((2^{n_1}+s(e_1)) \times 2^{n_2} + s(e_2))}$,
%
\item if  $\theta=\mbox{\prettyll{SInt}}$ and $s(e_1) < 0, s(e_2) < 0$ then \\  $s(cat(e_1, e_2))$ is $bv_{\mbox{\prettyll{UInt}}\bwidth{n_1+n_2}((2^{n_1}-s(e_1)) \times 2^{n_2} + 2^{n_2}-s(e_2))}$,
\end{itemize}
\end{itemize}
%
\item let $te(e_1) = \theta \bwidth{n_1}$ and $te(e_2) = \mbox{\prettyll{UInt}} \bwidth{n_2}$, where $\theta \in \{\mbox{\prettyll{UInt}}, \mbox{\prettyll{SInt}}\}$, then
\begin{itemize}
\item $te(dshl(e_1, e_2)) = \theta \bwidth{n_1+2^{n_2}-1}$ and $s(dshl(e_1, e_2)) = bv_{\theta \bwidth{n_1+2^{n_2}-1}(s(e_1) \times 2^{s(e_2)})}$,
%
\item $te(dshr(e_1, e_2)) = \theta \bwidth{n_1}$ and $s(dshr(e_1, e_2)) = bv_{\theta \bwidth{n_1}(s(e_1) \div 2^{s(e_2)})}$,
%
\end{itemize}
%
\item let $te(e_1) = \theta \bwidth{n_1}$, where $\theta \in \{\mbox{\prettyll{UInt}}, \mbox{\prettyll{SInt}}\}$, then
\begin{itemize}
\item $te(andr(e_1)) = te(orr(e_1)) =  te(xorr(e_1)) = \uint{1}$, moreover, $s(andr(e_1))$ (resp. $s(orr(e_1))$, $s(xorr(e_1))$) is the and (resp. or, xor) of all the bits in the bit vector $s(e_1)$,
%
\item $te(neg(e_1)) = \sint{n_1 +1}$, $te(not(e_1)) =  \uint{n_1}$, moreover,
\begin{itemize}
\item $s(neg(e_1)) = bv_{\mbox{\prettyll{SInt}}\bwidth{n_1+1}(-s(e_1))}$,
%
\item $s(not(e_1))$ is obtained by applying the logical not to each bit of the bit vector $s(e_1)$,
\end{itemize}
%\begin{itemize}
%\item if $\theta = UInt$, then $s(neg(e_1)) = bv_{SInt\bwidth{n_1+1}(-s(e_1))}$,
%
%\item if $\theta = SInt$, then $s(neg(e_1)) = bv_{SInt\bwidth{n_1+1}(-s(e_1))}$.
%\end{itemize}
%
\item $te(cvt(e_1)) = \sint{n_1 +1}$ if $\theta = \mbox{\prettyll{UInt}}$, and $te(cvt(e_1)) =  \sint{n_1}$ otherwise,  moreover,
\begin{itemize}
\item if $\theta = \mbox{\prettyll{UInt}}$, then $s(cvt(e_1)) = bv_{\mbox{\prettyll{SInt}}\bwidth{n_1+1}(s(e_1))}$,
%
\item if $\theta = \mbox{\prettyll{SInt}}$, then $s(cvt(e_1)) = s(e_1)$,
\end{itemize}
\end{itemize}
%
\item let $te(e_1) = \theta \bwidth{n_1}$, where $\theta \in \{\mbox{\prettyll{UInt}}, \mbox{\prettyll{SInt}}\}$, then
\begin{itemize}
\item $te(pad(e_1, n)) = \theta\bwidth{\max(n_1, n)}$, moreover, if $n_1 < n$, then $s(pad(e_1,n)) = bv_{\theta\bwidth{n}(s(e_1))}$, otherwise, $s(pad(e_1,n)) = s(e_1)$,
%
\item $te(shl(e_1,n)) = \theta \bwidth{n_1 + n}$, moreover, $s(shl(e_1,n)) =  bv_{\theta\bwidth{n_1 + n}(2^{n} \times s(e_1))}$,
%\begin{itemize}
%\item if $\theta = UInt$, then $s(shl(e_1,n)) =  bv_{UInt\bwidth{n_1 + n}(2^{n} \times s(e_1))}$,
%\item if $\theta = SInt$, then $s(shl(e_1,n)) = bv_{SInt\bwidth{n_1 + n}(2^{n} \times s(e_1))}$,
%\end{itemize}
%
\item $te(shr(e_1, n)) = \theta \bwidth{n_1-n}$ if $n_1 > n$, and $te(shr(e_1, n)) =  \theta \bwidth{1}$ otherwise,  moreover,
\begin{itemize}
\item if $n_1 > n$, then $s(shr(e_1, n))$ is obtained from $s(e_1)$ by shifting $n$ bits right,
%
\item otherwise, if $\theta = \mbox{\prettyll{UInt}}$, then $s(shr(e_1, n))=bv_{\mbox{\prettyll{UInt}}<1>(0)}$, otherwise, $s(shr(e_1, n))=bv_{\mbox{\prettyll{SInt}}<1>(0)}$ if $s(e_1) \ge 0$, and $s(shr(e_1, n))=bv_{\mbox{\prettyll{SInt}}<1>(1)}$ if $s(e_1) < 0$,
\end{itemize}
%
\item $te(head(e_1, n)) = { \mbox{\prettyll{UInt}}}{n}$, $te(tail(e_1, n)) = { \mbox{\prettyll{UInt}}}{n_1-n}$, moreover, $s(head(e_1, n))$ is the most significant $n$ bits of $s(e_1)$, $s(tail(e_1, n))$ is obtained from $s(e_1)$ by truncating the $n$ most significant bits of $s(e_1)$,
%
\item $te(bits(e_1, n, n')) = \uint{n'-n+1}$, and $s(bits(e_1, n, n'))$ is the bit vector comprising the $n, n+1, \cdots, n'$-th bit of $s(e_1)$.
\end{itemize}
%
\end{itemize}
}
%%%%%%%%%%%%%%
%%%%%%%%%%%%%%

Next, let us define the rules for $\updatecm$, $\updatecmref$, and $\invalidateref$.

The function $\updatecm(cm_1, id, ft, ls)$ is defined the following rule. It updates $cm_1$ into $cm_2$ by modifying $cm_1(v, m+n)$ to $ls[n]$ for each $0 \le n < sz$, where $id = (v, m)$, $sz = \sizeofftype(ft)$. 
%
\begin{gather*}
		\infer[]
		{\updatecm(cm_1, id, ft, ls) = cm_2}
		{
			\begin{array}{c}
				id = (v, m) \wedge \sizeofftype(ft) = sz \\
				\forall n.\ 0 \le n < sz \Rightarrow cm_2(v, m+n) = ls[n] \\
				\forall v', n'.\ (v' \neq v \vee n' < m \vee n' \ge m+sz) \Rightarrow cm_2(v', n') = cm_1(v', n') \\
			\end{array}
		}
\end{gather*}	

The function $\updatecmref(cm_1, id, ft, ls, id', ft', ls')$ is defined the following rule. It first calls $lgt = \listgtyperef(ft, false)$ to compute a list $lgt$ from $ft$ (the type of $r$), where each element in $lgt$ is a pair $(gt, flipped)$ such that $gt \in \gtypes$ and $flipped$ is a Boolean value denoting whether the corresponding field is flipped or not. Then for each $n: 0 \le n < sz$, it updates $cm_1(v, m+n)$ to $ls'[n]$ if $lgt[n] = (\_, \lfalse)$, and updates $cm_1(v', m'+n)$ to $ls[n]$ if $lg[n] = (\_, \ltrue)$, where $id = (v, m)$ and $id' = (v', m')$. 
%
\begin{gather*}
		\infer[]
		{\updatecmref(cm_1, id, ft, ls, id', ft', ls') = cm_2}
		{
			\begin{array}{c}
				id = (v, m) \wedge id' = (v', m') \hspace{15pt} \sizeofftype(ft) = sz \hspace{15pt} \listgtyperef(ft, \lfalse) = lgt \\
				\forall n.\ (0 \le n < sz \wedge lgt[n] = (\_, \lfalse)) \Rightarrow \\
				\big(
				%    \begin{array}{l}
					cm_2(v, m+n) = ls'[n] \wedge ((v',m') \neq (v, m)  \rightarrow cm_2(v', m' + n) = cm_1(v', m'+n)) 
					%    \end{array}
				\big)
				\\ 
				\forall n.\ (0 \le n < sz \wedge lgt[n] = (\_, \ltrue)) \Rightarrow \\
				\big(
				%    \begin{array}{l}
					cm_2(v', m'+n) = ls[n] \wedge  
					((v, m) \neq (v', m')  \rightarrow cm_2(v, m + n) = cm_1(v, m+n))
					%    \end{array}
				\big)
				\\
				\forall v'', n''.\ 
				\left(
				\begin{array}{l}
					(v'' \neq v \vee n'' < m \vee n'' \ge m+sz)\ \wedge \\
					(v'' \neq v' \vee n'' < m' \vee n'' \ge m'+sz)
				\end{array}
				\right)
				\Rightarrow cm_2(v'', n'') = cm_1(v'', n'') 
			\end{array}
		}
\end{gather*}

The function $\invalidateref(cm_1, id, ft, ori)$ is defined by the following rule. It first calls $lgt = \listgtyperef(ft, false)$ to compute a list $lgt$ from $ft$ (the type of $r$), where each element in $lgt$ is a pair $(gt, \flipped)$ such that $gt \in \gtypes$ and $\flipped$ is a Boolean value denoting whether the corresponding field is flipped or not. Then for each $n: 0 \le n < sz$, it updates $cm_1(v, m+n)$ to $\svinvalid$ if either ($lgt[n] = (\_, \lfalse)$ and $ori = \duplex$ or $\sink$), or ($lgt[n] = (\_, \ltrue)$ and $ori = \duplex$ or $\source$), where $id = (v, m)$. 
%
\begin{gather*}
		\infer[]
		{\invalidateref(cm_1, id, ft, ori) = cm_2}
		{
			\begin{array}{c}
				id = (v, m) \hspace{15pt} \sizeofftype(ft) = sz \hspace{15pt} \listgtyperef(ft, false) = lgt \\
				\forall n.\ 0 \le n < sz \wedge 
				\left(
				\begin{array}{l}
					(lgt[n] = (\_, false) \wedge (ori = \sink \vee ori = \duplex))\ \vee \\
					(lgt[n] = (\_, true) \wedge (ori = \source \vee ori = \duplex))
				\end{array}
				\right) \Rightarrow\\
				\hspace{1cm}  cm_2(v, m+n) = \svinvalid  \\
				\forall n.\ 0 \le n < sz \wedge 
				\left(
				\begin{array}{l}
					(lgt[n] = (\_, false) \wedge ori = \source)\ \vee \\
					(lgt[n] = (\_, true) \wedge ori = \sink)
				\end{array}
				\right)  
				\Rightarrow \\
				cm_2(v, m+n) = cm_1(v, m+n)  \\
				\forall v', n'.\ (v' \neq v \vee n' < m  \vee n' \ge m+sz)  \rightarrow cm_2(v', n') = cm_1(v', n')
			\end{array}
		}
\end{gather*}


%%%%%%%%%%% the original updatecm and updatecmref, invalidateref
\hide{
\begin{itemize}
	\item The function $\updatecm(cm_1, id, ft, ls)$ updates $cm_1$ into $cm_2$ by modifying $cm_1(v, m+n)$ to $ls[n]$ for each $0 \le n < sz$, where $id = (v, m)$, $sz = \sizeofftype(ft)$. 

	\item The function $\updatecmref(cm_1, id, ft, ls, id', ft', ls')$ first calls $lgt = \listgtyperef(ft, false)$ to compute a list $lgt$ from $ft$ (the type of $r$), where each element in $lgt$ is a pair $(gt, flipped)$ such that $gt \in \gtypes$ and $flipped$ is a Boolean value denoting whether the corresponding field is flipped or not. Then for each $n: 0 \le n < sz$, it updates $cm_1(v, m+n)$ to $ls'[n]$ if $lgt[n] = (\_, \lfalse)$, and updates $cm_1(v', m'+n)$ to $ls[n]$ if $lg[n] = (\_, \ltrue)$, where $id = (v, m)$ and $id' = (v', m')$. 
	%
	\item The function $\invalidateref(cm_1, id, ft, ori)$ first calls $lgt = \listgtyperef(ft, false)$ to compute a list $lgt$ from $ft$ (the type of $r$), where each element in $lgt$ is a pair $(gt, \flipped)$ such that $gt \in \gtypes$ and $\flipped$ is a Boolean value denoting whether the corresponding field is flipped or not. Then for each $n: 0 \le n < sz$, it updates $cm_1(v, m+n)$ to $\svinvalid$ if either ($lgt[n] = (\_, \lfalse)$ and $ori = \duplex$ or $\sink$), or ($lgt[n] = (\_, \ltrue)$ and $ori = \duplex$ or $\source$), where $id = (v, m)$. 
\end{itemize} 

%%%%%% updatecm, updatecmref
\begin{figure}[htbp]
	\begin{gather*}
		\infer[\updatecm]
		{\updatecm(cm_1, id, ft, ls) = cm_2}
		{
			\begin{array}{c}
				id = (v, m) \wedge \sizeofftype(ft) = sz \\
				\forall n.\ 0 \le n < sz \Rightarrow cm_2(v, m+n) = ls[n] \\
				\forall v', n'.\ (v' \neq v \vee n' < m \vee n' \ge m+sz) \Rightarrow cm_2(v', n') = cm_1(v', n') \\
			\end{array}
		}
		\\[1ex]
		%%%%%%%%%%%%%%%%%%%%%%%%%%%%%%%%%%%%%%%%%%%%%%%%%%%%%%%%%%%%%%%%%%%%%%%%%%%%%%%%%%%%%%%%%%%%%%  
		\infer[\updatecmref]
		{\updatecmref(cm_1, id, ft, ls, id', ft', ls') = cm_2}
		{
			\begin{array}{c}
				id = (v, m) \wedge id' = (v', m') \wedge \sizeofftype(ft) = sz\ \wedge \\
				\listgtyperef(ft, \lfalse) = lgt \\
				\forall n.\ (0 \le n < sz \wedge lgt[n] = (\_, \lfalse)) \Rightarrow \\
				\big(
				%    \begin{array}{l}
					cm_2(v, m+n) = ls'[n] \wedge ((v',m') \neq (v, m)  \rightarrow cm_2(v', m' + n) = cm_1(v', m'+n)) 
					%    \end{array}
				\big)
				\\ 
				\forall n.\ (0 \le n < sz \wedge lgt[n] = (\_, \ltrue)) \Rightarrow \\
				\big(
				%    \begin{array}{l}
					cm_2(v', m'+n) = ls[n] \wedge  
					((v, m) \neq (v', m')  \rightarrow cm_2(v, m + n) = cm_1(v, m+n))
					%    \end{array}
				\big)
				\\
				\forall v'', n''.\ 
				\left(
				\begin{array}{l}
					(v'' \neq v \vee n'' < m \vee n'' \ge m+sz)\ \wedge \\
					(v'' \neq v' \vee n'' < m' \vee n'' \ge m'+sz)
				\end{array}
				\right)
				\Rightarrow cm_2(v'', n'') = cm_1(v'', n'') 
			\end{array}
		}
		\\[1ex]
		\infer[\invalidateref]
		{\invalidateref(cm_1, id, ft, ori) = cm_2}
		{
			\begin{array}{c}
				id = (v, m) \wedge \sizeofftype(ft) = sz \wedge \listgtyperef(ft, false) = lgt \\
				\forall n.\ 0 \le n < sz \wedge 
				\left(
				\begin{array}{l}
					(lgt[n] = (\_, false) \wedge (ori = \sink \vee ori = \duplex))\ \vee \\
					(lgt[n] = (\_, true) \wedge (ori = \source \vee ori = \duplex))
				\end{array}
				\right) \Rightarrow\\
				\hspace{1cm}  cm_2(v, m+n) = \svinvalid  \\
				\forall n.\ 0 \le n < sz \wedge 
				\left(
				\begin{array}{l}
					(lgt[n] = (\_, false) \wedge ori = \source)\ \vee \\
					(lgt[n] = (\_, true) \wedge ori = \sink)
				\end{array}
				\right)  
				\Rightarrow \\
				cm_2(v, m+n) = cm_1(v, m+n)  \\
				\forall v', n'.\ (v' \neq v \vee n' < m  \vee n' \ge m+sz)  \rightarrow cm_2(v', n') = cm_1(v', n')
			\end{array}
		}
	\end{gather*}
	\caption{$\updatecm$, $\updatecmref$, and $\invalidateref$}
	\label{tab:update-vm-ct}
\end{figure}
}
%%%%%%%%%%%
%%%%%%%%%%%


The function $\listgexp(e, ft)$ is defined in Fig.~\ref{fig:list-gexp}. 
Let us explain the rules $\listgexp: mux$ and $\listgexp: vector$. The explanation of the other rules are omitted for brevity. 
%
\begin{itemize}
	\item In the rule $\listgexp: mux$, $e = \mux(c, e_1, e_2)$. Then it computes $l1 = \listgexp(e1, ft)$ and $l2 = \listgexp(e2, ft)$. Moreover, it checks that the lengths of $l1$ and $l2$ are equal. Finally, the result $l$ is computed as $\mux(c, l1, l2)$, where the operator $\mux(c, \_, \_)$ is applied pointwisely to the elements in $l1$ and $l2$. 
	%
	\item In the rule $\listgexp: vector$, $e = r[n]$. Then it computes $(ft', ori) = \typeofref(r)$ and $sz  = \sizeofftype(ft)$. The result $l$ is defined as $(l'[0])[n]\ ++\ \cdots ++\ (l'[sz-1])[n]$, where $l' = \listgexp(r, ft')$. 
\end{itemize}

%%%%%% listgexp
\begin{figure}[htbp]
	\begin{gather*}
		\infernone[\listgexp: ground]
		{\listgexp(e, (ground, t)) = [e]}
		%    {
			%    \begin{array}{c}
				%    	ft = (ground, t) \wedge l = [e]
				%    \end{array}
			%    }
		\quad
		\\[1ex]
		%%%%%%%%%%%%%%%%%%%%%%%%%%%%%%%%%%%%%%%%%%%%%%%%
		\infer[\listgexp: mux]
		{\listgexp(e, ft) = l}
		{
			\begin{array}{c}
				e = \mux(c, e1, e2) \hspace{15pt} \listgexp(e1, ft) = l1 \hspace{15pt} \listgexp(e2, ft) = l2\  \wedge\\
				length(l1) = length(l2) \hspace{15pt} l = \mux(c, l1, l2)
			\end{array}
		}
		\quad
		\\[1ex]
		\infer[\listgexp: vector]
		{\listgexp(e, ft) = l}
		{
			\begin{array}{c}
				e= r[n] \hspace{15pt} \typeofref(r) = (ft', ori) \hspace{15pt} \sizeofftype(ft) = sz \\
				l' = \listgexp(r, ft') \hspace{15pt} l = (l'[0])[n]\ ++ \ \cdots \ ++ (l'[sz-1])[n]
			\end{array}
		}
		\quad
		\\[1ex]
		\infernone[\listgexp: bundle]
		{\listgexp(e,  (bundle, ff) ) = \listgexpfields(e, ff)}
		%    {
			%    \begin{array}{c}
				%    	 ft = (bundle, ff) \wedge l = \listgexpfields(e, ff)
				%    \end{array}
			%    }
		\quad
		\\[2ex]
		\infernone[\listgexpfields: nil]
		{\listgexpfields(e, \nil) = []}
		%    {
			%    \begin{array}{c}
				%    	 ff = nil \wedge l = []
				%    \end{array}
			%    }
		%%%%%%%%%%%%%%%%%%%%%%%%%%%%%%%%%%%%%%%%%%%%%%%
		\quad
		\\[1ex]
		\infernone[\listgexpfields: $\neq$nil]
		{\listgexpfields(e, (f, b, t) ff1) = \listgexp(e.f, t)\ ++\ \listgexpfields(e, ff1)}
		%    {
			%    \begin{array}{c}
				%    	 ff = (f, b, t) ff1 \wedge l = \listgexp(e.f, t)\ ++\ \listgexpfields(e, ff1)
				%    \end{array}
			%    }
	\end{gather*}
	\caption{$\listgexp$ and $\listgexpfields$}
	\label{fig:list-gexp}
\end{figure}

The relations $\typecompatible(ft, ft', \flipped)$ and  $\typecompatiblefields(\myvec{f}, \myvec{f'}, \flipped)$ are defined by the following rules. 
\begin{itemize}
	\item In the rule ``exp-gt'', $ft$ and $ft'$ are explicit ground types, that is, types from $\{\uint, \sint\} \times \Nat \cup \{\clock, \reset, \asyncreset\}$, and required to be equal. 
	%
	\begin{gather*}
		\infer[exp-gt]
		{\typecompatible(ft, ft', \flipped)}
		{
			\begin{array}{c}
				ft = (\gtype, t) \hspace{15pt} ft' = (\gtype, t) \\
				t \in \{\uint, \sint\} \times \Nat \cup \{\clock, \reset, \asyncreset\}
			\end{array}
		}
	\end{gather*}
	\item In the rule ``imp-gt-1'',  $ft = (\gtype, t)$, $ft' = (\gtype, t')$, $t$ is implicit and $t'$ is explicit, that is, either $t= (\uintimp, n)$ and $t' = (\uint, n')$, or $t= (\sintimp, n)$ and $t' = (\sint, n')$.  If $\flipped = \lfalse$, then the connection flow is from $ft'$ to $ft$, thus it is required that $n \ge n'$. 
	\begin{gather*}
		\infer[imp-gt-1]
		{\typecompatible(ft, ft', \flipped)}
		{
			\begin{array}{c}
				ft = (\gtype, t) \hspace{15pt} ft' = (\gtype, t') \\
				(t = (\uintimp, n) \wedge t' = (\uint, n')) \vee (t = (\sintimp, n) \wedge t' = (\sint, n')) \\
				\flipped = \lfalse \rightarrow n \ge n'
			\end{array}
		}	
	\end{gather*}	
	%
	\item In the rule ``imp-gt-2'',  $ft = (\gtype, t)$, $ft' = (\gtype, t')$, $t$ is explicit and $t'$ is implicit, that is, either $t= (\uint, n)$ and $t' = (\uintimp, n')$, or $t= (\sint, n)$ and $t' = (\sintimp, n')$.  If $\flipped = \ltrue$, then the connection flow is from $ft$ to $ft'$, thus it is required $n \le n'$. 
	%
	\begin{gather*}
		\infer[imp-gt-2]
		{\typecompatible(ft, ft', \flipped)}
		{
			\begin{array}{c}
				ft = (\gtype, t) \hspace{15pt} ft' = (\gtype, t') \\
				(t = (\uint, n) \wedge t' = (\uintimp, n')) \vee (t = (\sint, n) \wedge t' = (\sintimp, n')) \\
				\flipped = \ltrue \rightarrow n \le n'
			\end{array}
		}	
	\end{gather*}
	\item In the rule ``imp-gt-3'',  $ft = (\gtype, t)$, $ft' = (\gtype, t')$, both $t$ and $t'$ are implicit, that is, $t= (\uintimp, n)$ and $t' = (\uintimp, n')$, or $t= (\sintimp, n)$ and $t' = (\sintimp, n')$.  If $\flipped = \lfalse$, then $n \ge n'$, and $n \le n'$ otherwise. 
	%
	\begin{gather*}
		\infer[imp-gt-3]
		{\typecompatible(ft, ft', \flipped)}
		{
			\begin{array}{c}
				ft = (\gtype, t) \hspace{15pt} ft' = (\gtype, t') \\
				(t = (\uintimp, n) \wedge t' = (\uintimp, n')) \vee (t = (\sintimp, n) \wedge t' = (\sintimp, n')) \\
				(\flipped = \lfalse \wedge n \ge n') \vee (\flipped = \ltrue \wedge n \le n')
			\end{array}
		}	
	\end{gather*}
	\item In the rule ``vector'', $ft = (\vtype, t[n])$ and $ft' = (\vtype, t'[n])$, then it is required that $t$ and $t'$ is compatible, that is, $\typecompatible(t, t', \flipped)$ holds. 
	%
	\begin{gather*}
		\infer[vector]
		{\typecompatible(ft, ft', \flipped)}
		{
			\begin{array}{c}
				ft = (\vtype, t[n]) \hspace{15pt} ft' = (\vtype, t'[n]) \hspace{15pt} \typecompatible(t, t', \flipped)
			\end{array}
		}	
	\end{gather*}
	\item In the rule ``bundle'', $ft = (\btype, \myvec{f})$ and $ft' = (\btype, \myvec{f'})$. Then it is required that $\typecompatiblefields(\myvec{f}, \myvec{f'}, \flipped)$ holds. 
	%
	\begin{gather*}
		\infer[bundle]
		{\typecompatible(ft, ft', \flipped)}
		{
		\begin{array}{c}
			ft = (\btype, \myvec{f}) \hspace{15pt} ft' = (\btype, \myvec{f'}) \hspace{15pt}  \typecompatiblefields(\myvec{f}, \myvec{f'}, \flipped)
		\end{array}
		}	
	\end{gather*}
	\item In the rule ``nil'', $\typecompatiblefields(\myvec{f}, \myvec{f'}, \flipped)$ holds iff $\myvec{f} = \myvec{f'} = \nil$. 
	%
	\begin{gather*}
		\infer[nil]
		{\typecompatiblefields(\myvec{f}, \myvec{f'}, \flipped)}
		{
			\begin{array}{c}
				\myvec{f} = \nil \hspace{15pt} \myvec{f'} = \nil
			\end{array}
		}	
	\end{gather*}
	\item In the rule ``$\neq$nil-false'', $\myvec{f} = (x, b, t)\ \myvec{f_1}$, $ff' = (x', b, t')\ \myvec{f'_1}$, and $b = \lfalse$. Then it is required that $\typecompatible(t, t', \flipped)$ and $\typecompatiblefields(\myvec{f_1}, \myvec{f'_1}, \flipped)$ hold. 
	%
	\begin{gather*}
		\infer[$\neq$nil-false]
		{\typecompatiblefields(\myvec{f}, \myvec{f'}, \flipped)}
		{
			\begin{array}{c}
				\myvec{f} = (x, b, t)\ \myvec{f_1} \hspace{15pt} \myvec{f'} = (x', b, t')\ \myvec{f'_1} \hspace{15pt} b = \lfalse \\
				\typecompatible(t, t', \flipped) \hspace{15pt}
				\typecompatiblefields(\myvec{f_1}, \myvec{f'_1}, \flipped)
			\end{array}
		}	
	\end{gather*}
	\item In the rule ``$\neq$nil-true'', $\myvec{f} = (x, b, t)\ \myvec{f_1}$, $\myvec{f'} = (x', b, t')\ \myvec{f'_1}$, and $b = \ltrue$. Then it is required that $\typecompatible(t, t', \neg \flipped)$ and $\typecompatiblefields(\myvec{f_1}, \myvec{f'_1}, \flipped)$  hold. Note that because $b = \ltrue$, that is, $x$ and $x'$ are flipped fields, the value $\neg \flipped$, instead of $\flipped$, is used in $\typecompatible(t, t', \neg \flipped)$. 
	\begin{gather*}
		\infer[$\neq$nil-true]
		{\typecompatiblefields(\myvec{f}, \myvec{f'}, \flipped)}
		{
			\begin{array}{c}
				\myvec{f} = (x, b, t)\ \myvec{f_1} \hspace{15pt} \myvec{f'} = (x', b, t') \myvec{f'_1} \hspace{15pt} b = \ltrue \\
				\typecompatible(t, t', \neg \flipped) \hspace{15pt}
				\typecompatiblefields(\myvec{f_1}, \myvec{f'_1}, \flipped)
			\end{array}
		}	
	\end{gather*}
\end{itemize}

%%%%%%%%%%%% the original definition of typecompatible and typecompatiblefields
%%%%%%%%%%%%
\hide{
%%%%%% typecompatible
\begin{figure}[htbp]
	\begin{gather*}
		\infer[exp-gt]
		{\typecompatible(ft, ft', \flipped)}
		{
			\begin{array}{c}
				ft = (\gtype, t) \wedge ft' = (\gtype, t) \\
				t \in \{\uint, \sint\} \times \Nat \cup \{\clock, \reset, \asyncreset\}
			\end{array}
		}
		\\[1ex]
		\infer[imp-gt-1]
		{\typecompatible(ft, ft', \flipped)}
		{
			\begin{array}{c}
				ft = (\gtype, t) \wedge ft' = (\gtype, t') \\
				(t = (\uintimp, n) \wedge t' = (\uint, n')) \vee (t = (\sintimp, n) \wedge t' = (\sint, n')) \\
				\flipped = \lfalse \rightarrow n \ge n'
			\end{array}
		}
		\\[1ex]
		\infer[imp-gt-2]
		{\typecompatible(ft, ft', \flipped)}
		{
			\begin{array}{c}
				ft = (\gtype, t) \wedge ft' = (\gtype, t') \\
				(t = (\uint, n) \wedge t' = (\uintimp, n')) \vee (t = (\sint, n) \wedge t' = (\sintimp, n')) \\
				\flipped = \ltrue \rightarrow n \le n'
			\end{array}
		}
		\\[1ex]
		\infer[imp-gt-3]
		{\typecompatible(ft, ft', \flipped)}
		{
			\begin{array}{c}
				ft = (\gtype, t) \wedge ft' = (\gtype, t') \\
				(t = (\uintimp, n) \wedge t' = (\uintimp, n')) \vee (t = (\sintimp, n) \wedge t' = (\sintimp, n')) \\
				(\flipped = \lfalse \wedge n \ge n') \vee (\flipped = \ltrue \wedge n \le n')
			\end{array}
		}
		\\[1ex]
		\infer[vector]
		{\typecompatible(ft, ft', \flipped)}
		{
			\begin{array}{c}
				ft = (\vtype, t[n]) \wedge ft' = (\vtype, t'[n]) \\
				\typecompatible(t, t', \flipped)
			\end{array}
		}
		\\[1ex]
		\inferii[bundle]
		{\typecompatible(ft, ft', \flipped)}
		{
			\seqq{ft = (\btype, \myvec{f}) \wedge ft' = (\btype, \myvec{f'})} 
		}
		{
			\seqq{\typecompatiblefields(\myvec{f}, \myvec{f'}, \flipped)}
		}
		\quad
		\\[1ex]
		\infer[nil]
		{\typecompatiblefields(\myvec{f}, \myvec{f'}, \flipped)}
		{
			\begin{array}{c}
				\myvec{f} = \nil \wedge \myvec{f'} = \nil
			\end{array}
		}
		\\[1ex]
		\infer[$\neq$nil-false]
		{\typecompatiblefields(\myvec{f}, \myvec{f'}, \flipped)}
		{
			\begin{array}{c}
				\myvec{f} = (x, b, t)\ \myvec{f_1} \wedge \myvec{f'} = (x', b, t')\ \myvec{f'_1} \wedge b = \lfalse \\
				\typecompatible(t, t', \flipped) \\
				\typecompatiblefields(\myvec{f_1}, \myvec{f'_1}, \flipped)
			\end{array}
		}
		\\[1ex]
		\infer[$\neq$nil-true]
		{\typecompatiblefields(\myvec{f}, \myvec{f'}, \flipped)}
		{
			\begin{array}{c}
				\myvec{f} = (x, b, t)\ \myvec{f_1} \wedge \myvec{f'} = (x', b, t') \myvec{f'_1} \wedge b = \ltrue \\
				\typecompatible(t, t', \neg \flipped) \\
				\typecompatiblefields(\myvec{f_1}, \myvec{f'_1}, \flipped)
			\end{array}
		}
	\end{gather*}
	\caption{$\typecompatible$ and $\typecompatiblefields$}
	\label{fig:type-compatible}
\end{figure}
}
%%%%%%%%%%%%
%%%%%%%%%%%%

The functions $\typeofref(r, tmap)$ and $\listgtyperef(ft, \flipped)$ compute the type and list of ground types respectively for references. 
Moreover, $\typeofref(r, tmap)$ calls $\typeoffield(f, \myvec{f'}, ori)$ to compute $(t, ori')$ such that $ori$ is the orientation of $\myvec{f'}$, $t$ is the type of $f$ in $\myvec{f'}$, and $ori'$ is the orientation of $f$. 
They are defined by the following rules.
% Fig.~\ref{fig:type-ref} and Fig.~\ref{fig:list-ftype-fields}.
%Let us explain some rules here. 
\begin{itemize}
	\item In the rule ``{\typeofref: id}'', if $r$ is a name $x$, then $\typeofref(r, tmap) $ is defined as $tmap(id(x))$. 
	\begin{gather*}
		\inferii[\typeofref: id]
		{\seqq{\typeofref(r, tmap) = (ft, ori)}}
		{\seqq{r = x}} {\seqq{tmap(id(x)) = (ft, ori)}}
	\end{gather*}
	%
	\item In the rule ``{\typeofref: vector}'', if $r= r'[n]$, then it recursively calls $\typeofref(r', tmap)$ to obtain $(t[m], ori)$. If $n  < m$, then it returns $(t, ori)$. 
	%
	\begin{gather*}
		\inferiii[\typeofref: vector]
		{\seqq{\typeofref(r, tmap) = (t, ori)}}
		{\seqq{r = r'[n]}} {\seqq{\typeofref(r', tmap) =  (t[m], ori)}} {\seqq{n < m}}
	\end{gather*}
	
	\item In the rule ``{\typeofref: field}'', if $r = r'.f$, then it computes $(ft', ori')= \typeofref(r', tmap)$ and returns $\typeoffield(f, ft', ori')$. 
	%
	\begin{gather*}
		\infer[\typeofref: field]
		{\seqq{\typeofref(r, tmap) = (ft, ori)}}
		{
		\begin{array}{c}
			\seqq{r = r'.f} \hspace{15pt} \typeofref(r', tmap) = (ft', ori') \\
			 \typeoffield(f, ft', ori') = (ft, ori)
		\end{array}
		}
	\end{gather*}
%
	\item In the rule ``{\typeoffield: nil}'', if $\myvec{f'} = \nil$, then it returns $\svnone$. 
	\begin{gather*}	
	\infer[\typeoffield: nil]
	{\seqq{\typeoffield(f, \myvec{f'}, ori) = \svnone}}
	{\seqq{\myvec{f'} = \nil}}
	\end{gather*}
%	
	\item In the rule ``\typeoffield: neq'', if $\myvec{f'} = (f_1, b, t)\ \myvec{f''}$ and  $f \neq f_1$, then it returns the result $\typeoffield(f, \myvec{f''}, ori)$. 
	\begin{gather*}
	\infer[\typeoffield: neq]
	{\seqq{\typeoffield(f, \myvec{f'}, ori) = (ft, ori')}}
	{
	\begin{array}{c}
		\myvec{f'} = (f_1, b, t)\ \myvec{f''} \hspace{15pt} f \neq f_1 \hspace{15pt}
		\typeoffield(f, \myvec{f''}, ori) = (ft, ori')
	\end{array}
	}	
	\end{gather*}
	\item In the rule ``{\typeoffield: eq-1}'', if $\myvec{f'} = (f_1, b, t)\ \myvec{f''}$ and $f = f_1$, moreover, $b = \lfalse$ or ($b = \ltrue$ and $ori = \duplex$), then  it returns $(t, ori)$. 
	%
	\begin{gather*}
	\infer[\typeoffield: eq-1]
	{\seqq{\typeoffield(f, \myvec{f'}, ori) = (t, ori)}}
	{
	\begin{array}{c}
		\myvec{f'} = (f_1, b, t)\ \myvec{f''} \hspace{15pt} f = f_1  \hspace{15pt}
		b = \lfalse \vee (b = \ltrue \wedge ori = \duplex)
	\end{array}
	}
	\end{gather*}

	\item In the rule ``{\typeoffield: eq-2}'', if $\myvec{f'} = (f_1, b, t)\ \myvec{f''}$ and $f = f_1$, then it returns $(t, \sink)$ if $b = \ltrue$ and $ori = \source$. 
	%
	\begin{gather*}
	\infer[\typeoffield: eq-2]
	{\seqq{\typeoffield(f, \myvec{f'}, ori) = (t, \sink)}}
	{
	\begin{array}{c}
		\myvec{f'} = (f_1, b, t)\ \myvec{f''} \hspace{15pt} f = f_1  \hspace{15pt}
		b = \ltrue \hspace{15pt} ori = \source
	\end{array}
	}
	\end{gather*}
%
	\item In the rule ``\typeoffield: eq-3'', if $\myvec{f'} = (f_1, b, t)\  \myvec{f''} $ and $f = f_1$, then it returns $(t, \source)$ if $b = \ltrue$ and $ori = \sink$. 
	\begin{gather*}
	\infer[\typeoffield: eq-3]
	{\seqq{\typeoffield(f, \myvec{f'}, ori) = (t, \source)}}
	{
	\begin{array}{c}
 		\myvec{f'} = (f_1, b, t)\  \myvec{f''} \hspace{15pt} f = f_1 \hspace{15pt} b = \ltrue \hspace{15pt} ori = \sink
	\end{array}
	}
	\end{gather*}
%	
	\item In the rule  ``\listgtyperef: ground'',  it simply returns $[(gt, \flipped)]$.
	
	\begin{gather*}
	\infer[\listgtyperef: ground]
	{
	 \seqq{\listgtyperef((\gtype, gt), \flipped)  = [(gt, \flipped)]}
	}
	{
	}
	\end{gather*}
%	
	\item In the rule ``{\listgtyperef: vector}'', if $ft = (\vtype, ft'[m])$ and $\listgtyperef(ft', \flipped) = l'$, then it returns $m$ copies of $l'$. %$\underbrace{l' \ ++\ \cdots\ ++\ l'}_{m}$.   
	%
	
	\begin{gather*}
	\infer[\listgtyperef: vector]
	{
	\seqq{\listgtyperef((\vtype, ft'[m]), \flipped) = \conc{m}(\listgtyperef(ft', \flipped))}
	}
	{ 
	}
	\end{gather*}
%
	\item In the rule ``{\listgtyperef: bundle}'', if $ft= (\btype, \myvec{f})$, then it returns 
	$\listgtypereffields(\myvec{f}, \flipped)$.
	%
	
	\begin{gather*}
		\infer[\listgtyperef: bundle]
		{
		 \listgtyperef(\btype, \myvec{f}, \flipped) =  \listgtypereffields(\myvec{f}, \flipped)
		}
		{
		}
	\end{gather*}
%	
	\item In the rule ``\listgtypereffields: nil'', it returns the empty list. 
	
	\begin{gather*}
		\infernone[\listgtypereffields: nil]
	{\seqq{\listgtypereffields(\nil, \flipped) = []}}
	\end{gather*}	
%
	\item In the rule ``{\listgtypereffields: $\neq$nil-1}'', if $\myvec{f} = (\_, \lfalse, ft')\  \myvec{f'}$, then it returns the concatenation of $\listgtyperef(ft', \flipped)$ and $\listgtypereffields(\myvec{f'}, \flipped)$. 
	\begin{gather*}
	\infer[\listgtypereffields: $\neq$nil-1]
	{\seqq{\listgtypereffields(\myvec{f}, \flipped) = l'\ {++}\ l''}}
	{
		\begin{array}{c}
			\myvec{f} = (\_, \lfalse, ft')\  \myvec{f'} \hspace{15pt} \listgtyperef(ft', \flipped) = l' \\
			 \listgtypereffields(\myvec{f'}, \flipped) = l''
		\end{array}
	}
	\end{gather*}	
%
	\item In the rule ``{\listgtypereffields: $\neq$nil-2}'', if   $\myvec{f} = (\_, \flipped', ft')\ \myvec{f'}$ and $\flipped' = \ltrue$, then it calls $\listgtyperef(ft', \neg \flipped)$ and $\listgtypereffields(\myvec{f'}, \flipped)$ to get $l'$ and $l''$ respectively, and finally returns $l' \ ++\ l''$. Note that $\neg \flipped$, instead of $\flipped$, is used in $\listgtyperef(ft', \neg \flipped)$ since $\flipped' = \ltrue$. 
	\begin{gather*}
	\infer[\listgtypereffields: $\neq$nil-2]
	{\seqq{\listgtypereffields(\myvec{f}, \flipped) = l'\ {++}\ l''}}
	{
		\begin{array}{c}
			\myvec{f} = (\_, \ltrue, ft')\  \myvec{f'} \hspace{15pt} \listgtyperef(ft', \neg \flipped) = l' \\
			 \listgtypereffields(\myvec{f'}, \flipped) = l''
		\end{array}
	}
	\end{gather*}
\end{itemize}


%%%%% removed 
%%%%% removed 
\hide{
%%%%%% typeofref
\begin{figure}[htbp]
	\def\defaultHypSeparation{\hskip 10pt}
	\begin{gather*}
		\inferii[\typeofref: id]
		{\seqq{\typeofref(r, tmap) = (ft, ori)}}
		{\seqq{r = x}} {\seqq{tmap(id(x)) = (ft, ori)}}
		\quad
		\\[1ex]
		\inferiii[\typeofref: vector]
		{\seqq{\typeofref(r, tmap) = (t, ori)}}
		{\seqq{r = r'[n]}} {\seqq{\typeofref(r', tmap) =  (t[m], ori)}} {\seqq{n < m}}
		\quad
		\\[1ex]
		%   \inferiii[\typeofref: vector-2]
		%    {\seqq{\typeofref(r, tmap) = \svnone}}
		%    {\seqq{r = r'[n]}} {\seqq{\typeofref(r', tmap) = (t[m], ori)}} {\seqq{n \ge m}}
		%    \quad
		%    \\[1ex]
		\def\defaultHypSeparation{\hskip 0pt}
		\inferiii[\typeofref: field]
		{\seqq{\typeofref(r, tmap) = (ft, ori)}}
		{
			\seqq{r = r'.f}
		}
		{
			\seqq{\typeofref(r', tmap) = (ft', ori')}
		}
		{
			\seqq{ \typeoffield(f, ft', ori') = (ft, ori)}
		}
		\quad
		%    \seqq{r = r'.f}} {\seqq{type\_of\_ref(r', tmap) = (ft', ori')}} {\seqq{type\_of\_field(f, ft', ori') = (ft, ori)}}
	\\[1ex]
	\infer[\typeoffield: \nil]
	{\seqq{\typeoffield(f, \myvec{f'}, ori) = \svnone}}
	{\seqq{\myvec{f'} = \nil}}
	\quad
	\\[1ex]
	\inferii[\typeoffield: neq]
	{\seqq{\typeoffield(f, \myvec{f'}, ori) = (ft, ori')}}
	{
		\seqq{\myvec{f'} = (f_1, b, t)\ \myvec{f''} \wedge f \neq f_1}
	}
	{
		\seqq{\typeoffield(f, \myvec{f''}, ori) = (ft, ori')}
	}
	\quad
	\\[1ex]
	\inferii[\typeoffield: eq-1]
	{\seqq{\typeoffield(f, \myvec{f'}, ori) = (t, ori)}}
	{
		\seqq{\myvec{f'} = (f_1, b, t)\ \myvec{f''} \wedge f = f_1}
	}
	{
		\seqq{b = \lfalse \vee (b = \ltrue \wedge ori = \duplex)}
	}
	\quad
	\\[1ex]
	\inferii[\typeoffield: eq-2]
	{\seqq{\typeoffield(f, \myvec{f'}, ori) = (t, \sink)}}
	{
		\seqq{\myvec{f'} = (f_1, b, t)\ \myvec{f''} \wedge f = f_1}
	}
	{
		\seqq{b = \ltrue \wedge ori = \source}
	}
	\quad
	\\[1ex]
	\inferii[\typeoffield: eq-3]
	{\seqq{\typeoffield(f, \myvec{f'}, ori) = (t, \source)}}
	{
		\seqq{
			\myvec{f'} = (f_1, b, t)\  \myvec{f''} \wedge f = f_1
		}
	}
	{
		\seqq{b = \ltrue \wedge ori = \sink}
	}
\end{gather*}
%  \vspace{-4mm}
\caption{The procedures $\typeofref$ and $\typeoffield$}
\label{fig:type-ref}
\end{figure}

\begin{align*}
\listgtyperef((\gtype, gt), \flipped) & = [(gt, \flipped)]\\ 
\listgtyperef((\vtype, ft'[m]), \flipped) & =  
\conc{m}(\listgtyperef(ft', \flipped))\\
\listgtyperef((\btype, \myvec{f}), \flipped) & = \listgtypereffields(\myvec{f}, \flipped)
\end{align*}

%%%%%% listgtyperef
\begin{figure}[htbp]
\begin{gather*}
	%    \infer[\listgtype: ground]
	%    {\seqq{list\_gtype(ft) = [gt]}}
	%    {\seqq{ft = (ground, gt)}}
	%    \quad
	%    \\[1ex]
	%    \inferii[list\_gtype: vector]
	%    {\seqq{list\_gtype(ft) = \underbrace{l'\ ++\ \cdots\ ++\ l'}_{m}}}
	%    {\seqq{ft = (vector, ft'[m])}}    {\seqq{l' = list\_gtype(ft')}}
	%    \quad
	%    \\[1ex]
	%    \inferii[list\_gtype: bundle]
	%    {\seqq{list\_gtype(ft) = l'}}
	%    {\seqq{ft = (bundle, ff)}}    {\seqq{l' = list\_gtype\_fields(ff)}}
	%    \quad
	%    \\[1ex]    
	%    \infer[list\_gtype\_fields: nil]
	%    {\seqq{list\_gtype\_fields(ff) = []}}
	%    {\seqq{f f = nil}}
	%    \quad
	%    \\[1ex]
	%    \inferiii[list\_gtype\_fields: $\neq$ nil]
	%    {\seqq{list\_gtype\_fields(ff) = l'\ {++}\ l''}}
	%   {\seqq{f f = (\_, \_, ft', ff')}} 
	%    {\seqq{list\_gtype(ft') = l'}} {\seqq{list\_gtype\_fields(ff') = l''}}  
	%    \quad
	%    \\[1ex]
	%    \infer[\listgtyperef: ground]
	%    {\seqq{\listgtyperef(ft, flipped) = [(gt, flipped)]}}
	%    {\seqq{ft = (\gtype, gt)}}
	%    \quad
	%    \\[1ex]
	%    \inferii[\listgtyperef: vector]
	%    {\seqq{\listgtyperef(ft, flipped) = \underbrace{l'\ ++\ \cdots\ ++\ l'}_{m}}}
	%    {\seqq{ft = (\vtype, ft'[m])}}    {\seqq{l' = \listgtyperef(ft', flipped)}}
	%    \quad
	%    \\[1ex]
	%    \inferii[\listgtyperef: bundle]
	%    {\seqq{\listgtyperef(ft, flipped) = l'}}
	%    {\seqq{ft = (\btype, \myvec{f})}}    {\seqq{l' = \listgtypereffields(\myvec{f}, flipped)}}
	%    \quad
	%    \\[1ex]
	%%%%%%%%%%%%%%%%%%%%%%%%%%%%%%%%%%%%%%%%%%%%%%%%%%%%%%%%%%%%%%%%%%%%%%%%%%%%%%%%%%%%%%%%%%%%%%     
	\infernone[\listgtypereffields: nil]
	{\seqq{\listgtypereffields(\nil, \flipped) = []}}
	%{\seqq{\myvec{f} = \nil}}
	\quad
	\\[1ex]
	\inferii[\listgtypereffields: $\neq$nil-1]
	{\seqq{\listgtypereffields(\myvec{f}, \flipped) = l'\ {++}\ l''}}
	{
		\begin{array}{c}
			\myvec{f} = (\_, \lfalse, ft')\  \myvec{f'} \ \wedge
		\end{array}
	}
	{
		\begin{array}{c}
			\listgtyperef(ft', \flipped) = l'\ \wedge \\
			\listgtypereffields(\myvec{f'}, \flipped) = l''
		\end{array}
	} 
	\quad
	\\[1ex]
	\inferii[\listgtypereffields: $\neq$nil-2]
	{\seqq{\listgtypereffields(\myvec{f}, \flipped) = l'\ {++}\ l''}}
	{
		\begin{array}{c}
			\myvec{f} = (\_, \ltrue, ft') \myvec{f'}\ \wedge 
		\end{array}
	}
	{
		\begin{array}{c}
			\listgtyperef(ft', \neg \flipped) = l'\ \wedge \\
			\listgtypereffields(\myvec{f'}, \flipped) = l''
		\end{array}
	} 
\end{gather*}
%  \vspace{-4mm}
\caption{$\listgtyperef$ and $\listgtypereffields$}
\label{fig:list-ftype-fields}
\end{figure}
}






\section{Missing details in Section~\ref{sec-form-verif-transform}}\label{app-hifirrtl}


%!TEX root = main.tex

\begin{figure}[htbp]
  \begin{gather*}
     \inferii[ ]
    {\seqq{{\sf portsStmtsTmap'}(pp, ss) = (tmap', var2exprs)}}
    {\seqq{{\sf portsTmap'}(pp) = pmap'}}
    {\seqq{{\sf stmtsTmap'}(pmap', ss) = (tmap', var2exprs)}}
   \quad
   \\[1ex]
    \infer[{\sf portsTmap'}: {\sf input}]
    {\seqq{{\sf portsTmap'}(pp) = pmap'}}
    {
    \begin{array}{c}
    p p= {\sf input}\ x: t\ pp' \wedge {\sf portsTmap'}(pp') = pmap \wedge pmap(id(x)) = \svnone \\
    pmap' = \fadd((id(x), (t, \source)), pmap)
    \end{array}
    }
    \quad
    \\[1ex]
    \infer[{\sf portsTmap'}: {\sf output}]
    {\seqq{{\sf portsTmap'}(pp) = pmap'}}
    {
    \begin{array}{c}
    p p= {\sf output}\ x: t\ pp' \wedge {\sf portsTmap'}(pp') = pmap \wedge pmap(id(x)) = \svnone \\
    pmap' = \fadd((id(x), (t, \sink)), pmap)
    \end{array}
    }
    \quad
    \\[1ex]
    \inferii[{\sf stmtsTmap'}: \nil]
    {\seqq{{\sf stmtsTmap'}((tmap, var2exprs), ss) = (tmap', var2exprs')}}
    {\seqq{s s = \nil}}
    {\seqq{tmap' = tmap \wedge var2exprs' = var2exprs}}
    \quad
    \\[1ex]
    \infer[{\sf stmtsTmap'}: $\neq$ \nil]
    {\seqq{{\sf stmtsTmap'}((tmap, var2exprs), ss) = (tmap', var2exprs')}}
    {
    \begin{array}{c}
      ss = s\ ss' \\
      (tmap'', var2exprs'') = stmtTmap((tmap, var2exprs), s) \\
      (tmap', var2exprs') = stmtsTmap((tmap'', var2exprs''), ss')
    \end{array}
    }
  \end{gather*}
%  \vspace{-4mm}
  \caption{${\sf portsStmtsTmap'}$, ${\sf portsTmap'}$, and ${\sf stmtsTmap'}$}
  \label{fig:ports_stmts_tmap'}
\end{figure}

\begin{figure}[htbp]
  \begin{gather*}
   \infer[find-some]
    {\seqq{{\sf addExprConnect}(lhsl, rhsl, var2exprs) = var2exprs'}}
    {
    \begin{array}{c}
    lhsl = (ref1, \_, \_) \ ltl \wedge rhsl = ehd \ etl \\
    var2exprs(ref1) = el \wedge var2exprs'' = \fadd(ref1, ehd :: el, var2exprs) \\
    var2exprs' = {\sf addExprConnect}(ltl, etl, var2exprs'')
    \end{array}
    }
    \quad
    \\[1ex]
    \infer[find-none]
    {\seqq{{\sf addExprConnect}(lhsl, rhsl, var2exprs) = var2exprs'}}
    {
    \begin{array}{c}
    lhsl = (ref1, \_, \_) \ ltl \wedge rhsl = ehd \ etl \\
    var2exprs(ref1) = \svnone \wedge var2exprs'' = \fadd(ref1, [::ehd], var2exprs) \\
    var2exprs' = {\sf addExprConnect}(ltl, etl, var2exprs'')
    \end{array}
    }
     \quad
    \\[1ex]
    \infer[flip, find-none]
    {\seqq{{\sf addExprConnectNonPassive}(lhsl, rhsl, var2exprs) = var2exprs'}}
    {
    \begin{array}{c}
    lhsl = (ref1, \_, true) \ ltl \wedge rhsl = (ref2, \_, \_) \ rtl \\
    var2exprs(ref2) = \svnone \wedge var2exprs'' = \fadd(ref2, [::(Eref \ ref1)], var2exprs) \\
    var2exprs' = {\sf addExprConnectNonPassive}(ltl, rtl, var2exprs'')
    \end{array}
    }
     \quad
    \\[1ex]
    \infer[non-flip, find-none]
    {\seqq{{\sf addExprConnectNonPassive}(lhsl, rhsl, var2exprs) = var2exprs'}}
    {
    \begin{array}{c}
    lhsl = (ref1, \_, false) \ ltl \wedge rhsl = (ref2, \_, \_) \ rtl \\
    var2exprs(ref1) = \svnone \wedge var2exprs'' = \fadd(ref1, [::(Eref \ ref2)], var2exprs) \\
    var2exprs' = {\sf addExprConnectNonPassive}(ltl, rtl, var2exprs'')
    \end{array}
    }
    \quad
    \\[1ex]
    \infer[flip, find-some]
    {\seqq{{\sf addExprConnectNonPassive}(lhsl, rhsl, var2exprs) = var2exprs'}}
    {
    \begin{array}{c}
    lhsl = (ref1, \_, true) \ ltl \wedge rhsl = (ref2, \_, \_) \ rtl \\
    var2exprs(ref2) = el \wedge var2exprs'' = \fadd(ref2, (Eref \ ref1) :: el, var2exprs) \\
    var2exprs' = {\sf addExprConnectNonPassive}(ltl, rtl, var2exprs'')
    \end{array}
    }
    \quad
    \\[1ex]
    \infer[non-flip, find-some]
    {\seqq{{\sf addExprConnectNonPassive}(lhsl, rhsl, var2exprs) = var2exprs'}}
    {
    \begin{array}{c}
    lhsl = (ref1, \_, false) \ ltl \wedge rhsl = (ref2, \_, \_) \ rtl \\
    var2exprs(ref1) = el \wedge var2exprs'' = \fadd(ref1, (Eref \ ref2) :: el, var2exprs) \\
    var2exprs' = {\sf addExprConnectNonPassive}(ltl, rtl, var2exprs'')
    \end{array}
    }
  \end{gather*}
\caption{${\sf addExprConnectNonPassive}$ and ${\sf addExprConnect}$}
  \label{tab:add_expr_connect}
\end{figure}


%%%%% where the bug begins
\begin{figure}[htbp]
   \begin{gather*}
    \infer[{\sf stmtTmap'}: skip]
    {\seqq{{\sf stmtTmap'}((tmap, var2exprs), s)=(tmap, var2exprs)}}
    {\seqq{s = skip}}
   \quad
   \\[1ex]
    \infer[{\sf stmtTmap'}: connect-ref]
    {\seqq{{\sf stmtTmap'}((tmap, var2exprs), s) = (tmap, var2exprs')}}
    {
    \begin{array}{c}
      s = r <= ref \\
      {\sf typeOfRef'}(r, tmap) = (t\_tgt, flip\_tgt)  \\
      {\sf typeOfRef'}(ref, tmap) = (t\_src, flip\_src) \\ 
      lhsl = {\sf listGRef} (r, t\_tgt, flip\_tgt) \\
      rhsl = {\sf listGRef} (ref, t\_src, flip\_src) \\
      var2exprs' = {\sf addExprConnectNonPassive} ((lhsl, rhsl), var2exprs) 
     \end{array}
    }
   \quad
   \\[1ex]
    \infer[{\sf stmtTmap'}: connect-expr]
    {\seqq{{\sf stmtTmap'}((tmap, var2exprs), s) = (tmap, var2exprs')}}
    {
    \begin{array}{c}
      s = r <= exp \\
      {\sf typeOfRef'}(r, tmap) = (t\_tgt, \_)  \\
      {\sf typeOfExp}(exp, tmap) = t\_src \\ 
      lhsl = {\sf listGRef} (r, t\_tgt, \lfalse) \\
      rhsl = {\sf listGExp} (exp, t\_src)\\
      var2exprs' = {\sf addExprConnect} ((lhsl, rhsl), var2exprs)
     \end{array}
    }
   \quad
   \\[1ex]
    \infer[{\sf stmtTmap'}: invalid]
    {\seqq{{\sf stmtTmap'}((tmap, var2exprs), s) = (tmap, var2exprs)}}
    {
      \begin{array}{c}
      s = r \mbox{ is invalid} \\
      {\sf typeOfRef'}(r, tmap) \neq \svnone
     \end{array}
    }
   \quad
   \\[1ex]
    \infer[{\sf stmtTmap'}: wire]
    {\seqq{{\sf stmtTmap'}((tmap, var2exprs), s) = (tmap', var2exprs)}}
    {
      \begin{array}{c}
      s = \mbox{wire } x: t \wedge tmap(id(x)) = \svnone  \\
      tmap' = \fadd((id(x), (t, \duplex)), tmap) 
     \end{array}
    }
   \quad
   \\[1ex]
    \infer[{\sf stmtTmap'}: node]
    {\seqq{{\sf stmtTmap'}((tmap, var2exprs), s) = (tmap', var2exprs')}}
    {
     \begin{array}{c}
      s = \mbox{node } x = e \wedge tmap(id(x)) = \svnone  \\
      {\sf typeOfExp}(e, tmap) = t\_src \\
      tmap' = \fadd((id(x), (t\_src, \source)), tmap) \\
      lhsl = {\sf listGRef} (x, t\_src, false) \\
      rhsl = {\sf listGExp} (e, t\_src)\\
      var2exprs' = {\sf addExprConnect} ((lhsl, rhsl),var2exprs )
     \end{array}
    }
\end{gather*}
%  \vspace{-4mm}
  \caption{{\sf stmtTmap'} (part 1)}
  \label{fig:stmts-tmap'}
\end{figure}

\begin{figure}[htbp]
  \small
   \begin{gather*}
    \infer[{\sf stmtTmap'}: register-reset]
    {\seqq{{\sf stmtTmap'}((tmap, var2exprs), s) = (tmap', var2exprs')}}
    {
    \begin{array}{c} 
      s = \mbox{reg } x: t \ clk \mbox{ with: reset =>} (rsig, rval)  \\
      tmap(id(x)) = \svnone \wedge {\sf typeOfExp}(clk, tmap) \neq \svnone  \\
      tmap' = \fadd((id(x), (t, \duplex)), tmap) \\
      lhsl = {\sf listGRef} (x, t, false) \\
      rhsl = {\sf listGExp} (rval, t)\\
      var2exprs' = {\sf addExprConnect} ((lhsl, rhsl), var2exprs )
     \end{array}
    }
   \quad
   \\[1ex]
    \infer[{\sf stmtTmap'}: register-non-reset]
    {\seqq{{\sf stmtTmap'}((tmap, var2exprs), s) = (tmap', var2exprs)}}
    {
    \begin{array}{c}
      s = \mbox{reg } x: t \ clk \\
      tmap(id(x)) = \svnone \wedge {\sf typeOfExp}(clk, tmap) \neq \svnone  \\
      tmap' = \fadd((id(x), (t, \duplex)), tmap) 
     \end{array}
    }
   \quad
   \\[1ex]
   \infer[{\sf stmtTmap'}: when]
    {\seqq{{\sf stmtTmap'}((tmap, var2exprs), s) = (tmap', var2exprs')}}
    {
    \begin{array}{c}
      s = \mbox{when } c: sst \mbox{ else }: ssf  \\
      {\sf typeOfExp}(c, scope) \neq \svnone \\
      {\sf stmtsTmap'}((tmap, var2exprs), sst) = (tmap'', var2exprs'') \\
      {\sf stmtsTmap'}((tmap'', var2exprs''), ssf) = (tmap', var2exprs')
     \end{array}
    }
    \end{gather*}
%  \vspace{-4mm}
  \caption{${\sf stmtTmap'} (part 2)$}
  \label{tab:stmts-tmap-2}
\end{figure}

\begin{figure}[htbp]
  \small
  \begin{gather*}
   \inferii[{\sf typeOfRef'}: id]
    {\seqq{{\sf typeOfRef'}(r, tmap) = (ft, false)}}
    {\seqq{r = x}} {\seqq{tmap(id(x)) = (ft, \_)}}
    \quad
    \\[1ex]
   \inferiii[{\sf typeOfRef'}: vector-1]
    {\seqq{{\sf typeOfRef'}(r, tmap) = (t, flip)}}
    {\seqq{r = r'[n]}} {\seqq{{\sf typeOfRef'}(r', tmap) =  (t[m], flip)}} {\seqq{n < m}}
    \quad
    \\[1ex]
   \inferiii[{\sf typeOfRef'}: vector-2]
    {\seqq{{\sf typeOfRef'}(r, tmap) = \svnone}}
    {\seqq{r = r'[n]}} {\seqq{{\sf typeOfRef'}(r', tmap) = (t[m], flip)}} {\seqq{n \ge m}}
    \quad
    \\[1ex]
   \infer[{\sf typeOfRef'}: field]
    {\seqq{{\sf typeOfRef'}(r, tmap) = {\sf typeOfField'}(f, \myvec{ff'}, flip')}}
    {
    \begin{array}{c}
    r = r'.f \\
    {\sf typeOfRef'}(r', tmap) = ((\btype, \myvec{ff'}), flip')
    \end{array}
    }
    \quad
    \\[1ex]
   \infer[{\sf typeOfField'} : \nil]
    {\seqq{{\sf typeOfField'}(f, \myvec{ff}, flip) = \svnone}}
    {\seqq{\myvec{ff} = \nil}}
    \quad
    \\[1ex]
   \infer[{\sf typeOfField'}: neq]
    {\seqq{{\sf typeOfField'}(f, \myvec{ff}, flip) = {\sf typeOfField'}(f, \myvec{ff'}, flip)}}
    {
    \begin{array}{c}
    \myvec{ff} = (f', flip, t, \myvec{ff'})\wedge f \neq f'
    \end{array}
    }
    \quad
    \\[1ex]
   \infer[{\sf typeOfField'}: eq-1]
    {\seqq{{\sf typeOfField'}(f, \myvec{ff}, flip) = (t, flip)}}
    {
    \begin{array}{c}
    \myvec{ff} = (f', flip, t, \myvec{ff'}) \wedge f = f' \\
    flipped = false
    \end{array}
    }
    \quad
    \\[1ex]
   \infer[{\sf typeOfField'}: eq-2]
    {\seqq{{\sf typeOfField'}(f, \myvec{ff}, flip) = (t, \neg flip)}}
    {
    \begin{array}{c}
    \myvec{ff} = (f', flip, t, \myvec{ff'}) \wedge f = f' \\
    flip = true
    \end{array}
    }
  \end{gather*}
%  \vspace{-4mm}
  \caption{${\sf typeOfRef'}$ and ${\sf typeOfField'}$}
  \label{tab:type-ref'}
\end{figure}

\begin{figure}[htbp]
  \begin{gather*}
    \inferii[{\sf ftSetSub}: ground-0]
    {\seqq{{\sf ftSetSub}(gt, n, ft) = [(\gtype, gt)]}}
    {\seqq{ft = (\gtype, \_)}}
    {\seqq{n = 0}}
    \quad
    \\[1ex]
    \inferii[{\sf ftSetSub}: ground-$\neq$0]
    {\seqq{{\sf ftSetSub}(gt, n, ft) = \svnone}}
    {\seqq{ft = (\gtype, \_)}}
    {\seqq{n \neq 0}}
    \quad
    \\[1ex]
    \infer[{\sf ftSetSub}: vector-lt]
    {\seqq{{\sf ftSetSub}(gt,n,ft) = (\vtype, newt[m])}}
    {
    \begin{array}{c}
    ft = (\vtype, ft'[m]) \wedge n < \sizeofftype(ft) \\
    n' = {\sf mod}(n, ({\sizeofftype(ft')})) \wedge newt = {\sf ftSetSub}(gt,n',ft')
    \end{array}
    }
    \quad
    \\[1ex]
    \infer[{\sf ftSetSub}: vector-geq]
    {\seqq{{\sf ftSetSub}(gt,n,ft) = \svnone}}
    {
    \begin{array}{c}
    ft = (\vtype, ft'[m]) \wedge n \geq \sizeofftype(ft) 
    \end{array}
    }
    \quad
    \\[1ex]
    \infer[{\sf ftSetSub}: bundle]
    {\seqq{{\sf ftSetSub}(gt,n,ft) = (\btype, \myvec{f'})}}
    {
    \begin{array}{c}
    ft = (\btype, \myvec{f}) \wedge \myvec{f'} = {\sf ftSetFld}(gt,n,\myvec{f})
    \end{array}
    }
    \quad
    \\[1ex]    
    \infer[{\sf ftSetFld}: nil]
    {\seqq{{\sf ftSetFld}(gt,n,\myvec{f}) = \svnone}}
    {\seqq{\myvec{f} = \nil}}
    \quad
    \\[1ex]
    \infer[{\sf ftSetFld}: $\neq$nil-1]
    {\seqq{{\sf ftSetFld}(gt,n,\myvec{f}) = (v, flip, ft', \myvec{f'})}}
    {
    \begin{array}{c}
    \myvec{f} = (v, flip, ft, \myvec{f'}) \ \wedge n < \sizeofftype(ft) \\
	ft' = {\sf ftSetSub}(gt,n,ft)
    \end{array}
    }
    \quad
    \\[1ex]
    \infer[{\sf ftSetFld}: $\neq$nil-2]
    {\seqq{{\sf ftSetFld}(gt,n,\myvec{f}) = (v, flip, ft, \myvec{f''})}}
    {
    \begin{array}{c}
    \myvec{f} = (v, flip, ft, \myvec{f'})\ \wedge \sizeofftype(ft) \leq n < \sizeoffields(\myvec{f}) \\
	\myvec{f''} = {\sf ftSetFld}(gt,(n- \sizeofftype(ft)),\myvec{f'})
    \end{array}
    }
    \quad
    \\[1ex]
    \infer[{\sf ftSetFld}: $\neq$nil-3]
    {\seqq{{\sf ftSetFld}(gt,n,\myvec{f}) = \svnone}}
    {
    \begin{array}{c}
    \myvec{f} = (v, flip, ft, \myvec{f'}) \ \wedge n \geq \sizeoffields(f) 
    \end{array}
    }
    \end{gather*}
    \caption{${\sf ftSetSub}$ and ${\sf ftSetFld}$}
  \label{tab:ft-set-sub}
\end{figure}

\begin{table}[htbp]
  \small
  \begin{gather*}
    \infer[{\sf drawG}: \nil]
    {\seqq{{\sf drawG}({\sf elements} (var2exprs), tmap) = (\nil, emptyg)}}
    {\seqq{{\sf elements} (var2exprs) = \nil}}
    \quad
   \\[1ex]
    \infer[{\sf drawG}: $\neq$ \nil, explicit]
    {\seqq{{\sf drawG}({\sf elements} (var2exprs), tmap, vertices, g) = {\sf drawG} (tl, tmap, vertices, g)}}
    {
    \begin{array}{c}
      {\sf elements} (var2exprs) = v : \_\ tl \\
      {\sf typeOfRef'}(v, tmap) = (t, \_) \wedge {\sf notImplict} (t) = \ltrue
    \end{array}
    }
    \quad
   \\[1ex]
    \infer[{\sf drawG}: $\neq$ \nil, implicit]
    {\seqq{{\sf drawG}({\sf elements} (var2exprs), tmap, vertices, g) = {\sf drawG} (tl, tmap, vertices'', g'')}}
    {
    \begin{array}{c}
      {\sf elements} (var2exprs) = v : el\ tl \\
      {\sf typeOfRef'}(v, tmap) = (t, \_) \wedge {\sf notImplict} (t) = false \\
      (vertices'', g'') = {\sf drawEl}(v, el, vertices, g) 
    \end{array}
    }
    \quad
   \\[1ex]
    \infer[{\sf drawEl}: \nil]
    {\seqq{{\sf drawEl}(v, el, vertices, g) = (vertices, g)}}
    {\seqq{el = \nil}}
    \quad
   \\[1ex]
    \infer[{\sf drawEl}: $\neq$ \nil]
    {\seqq{{\sf drawEl}(v, el, vertices, g) = {\sf drawEl}(v, tl, vertices', g')}}
    {
    \begin{array}{c}
      el = e\ tl \\
      vl = {\sf expr2varlist} (e) \wedge vertices' = vertices ++ vl \\
      g' = {\sf foldLeft}({\sf fun} (tv, tg) => \fadd(v, tv :: tg(v), tg)) \leftarrow (vl,g)
    \end{array}
    }
    \quad
    \\[1ex]
    \infer[{\sf expr2varlist}: const]
    {\seqq{{\sf expr2varlist}(e) = \nil}}
    {\seqq{e = const}}
   \quad
   \\[1ex]
    \infer[{\sf expr2varlist}: ref]
    {\seqq{{\sf expr2varlist}(e) = [::ref]}}
    {\seqq{e = ref}}
   \quad
   \\[1ex]
    \inferii[{\sf expr2varlist}: cast]
    {\seqq{{\sf expr2varlist}(e) = vl}}
    {\seqq{e = cast(e0)}}
    {\seqq{vl = {\sf expr2varlist}(e0)}}
   \quad
   \\[1ex]
    \inferii[{\sf expr2varlist}: unop]
    {\seqq{{\sf expr2varlist}(e) = vl}}
    {\seqq{e = unop(e0)}}
    {\seqq{vl = {\sf expr2varlist}(e0)}}
   \quad
   \\[1ex]
    \inferii[{\sf expr2varlist}: binop]
    {\seqq{{\sf expr2varlist}(e) = vl}}
    {\seqq{e = binop(e0, e1)}}
    {\seqq{vl = {\sf expr2varlist}(e0) ++ {\sf expr2varlist}(e1)}}
    \quad
   \\[1ex]
    \inferii[{\sf expr2varlist}: mux]
    {\seqq{{\sf expr2varlist}(e) = vl}}
    {\seqq{e = mux(e0, e1, e2)}}
    {\seqq{vl = {\sf expr2varlist}(e0) ++ {\sf expr2varlist}(e1) ++ {\sf expr2varlist}(e2)}}
  \end{gather*}
%  \vspace{-4mm}
  \caption{${\sf drawG}$, ${\sf drawEl}$, and ${\sf expr2varlist}$}
  \label{tab:ports-stmts-tmap}
  \vspace{-6mm}
\end{table}



%%%%%%%%%%%%%%%%%%%%%%%%%%%%%%%
%%%%%%%%%%%%%%%%%%%%%%%%%%%%%%%
\hide{
\section{Axiomatic semantics of compilation passes}\label{sec-appen-sem-cpass}

The semantics of compilation pass InferType on a module $M$ is specified by the following formula,
\[
\begin{array}{l c l}
\varphi_{InferType, M}(ce, cm, ce', cm') & \egdef & \exists (ce_i)_{0 \le i \le n} (cm_i)_{0 \le i \le n}.\ ce = ce_0 \wedge cm = cm_0 \wedge ce' = ce_n\ \wedge \\
& &  cm' = cm_n \wedge \bigwedge \limits_{1 \le i \le n} \varphi_{InferType, s_i}(ce_{i-1}, cm_{i-1}, ce_i, cm_i),
\end{array}
\]
where the semantics of $\varphi_{InferType, s_i}$ is defined by the rules in Table~\ref{tab:infertype-sem}. Moreover,
$type_{ce}(e)$ in Table~\ref{tab:infertype-sem} is defined in the same flavor as $te(e)$ in the definition of operational semantics of LoFIRRTL. For instance, $type_{ce}(add(e_1, e_2))$ is defined as follows:
\begin{itemize}
\item if $type_{ce}(e_1) = (\mbox{``Gtype''}, \theta\mbox{<}n_1\mbox{>})$ and  $type_{ce}(e_2) = (\mbox{``Gtype''}, \theta\mbox{<}n_2\mbox{>})$ for some $\theta \in \{\mbox{\prettyll{UInt}}, \mbox{\prettyll{SInt}}\}$, then $type_{ce}(add(e_1, e_2)) = (\mbox{``Gtype''}, \theta\mbox{<}\max(n_1,n_2)\mbox{>})$,
%\item if $type_{ce}(e_1) = (\mbox{``Gtype''}, \sint{n_1})$ and  $type_{ce}(e_2) = (\mbox{``Gtype''}, \sint{n_2})$, then $type_{ce}(add(e_1, e_2)) = (\mbox{``Gtype''}, \sint{\max(n_1,n_2)})$,
\item otherwise, $type_{ce}(add(e_1, e_2)) = \bot$.
\end{itemize}

\begin{table}[htbp]
  \footnotesize
  \begin{gather*}
%    \infer[skip]
%    {\seqq{\varphi_{InferType, \mbox{skip}}((ce, cm), (ce, cm))}}
%    {\seqq{-}}
%    \quad
    \infer[wire]
%    {\seqq{}}
    {\seqq{\varphi_{InferType, \mbox{wire }v:t}((ce, cm), (ce', cm))}}
   {\seqq{ce' = ce [(\mbox{``AggrType''}, t, \mbox{``Wire''})/v]}}
    \\[1ex]
    \infer[reg]
    {\varphi_{InferType, \mbox{reg } r:t, c \mbox{ with: reset => }(e_1, e_2)}((ce, cm), (ce', cm))}
   {\seqq{ce' = ce [(\mbox{``RegType''}, t, \mbox{``Register''})/r]}}
    \quad
    \\[1ex]
    \infer[node]
    {\seqq{\varphi_{InferType, \mbox{node } v = e} ((ce, cm), (ce', cm'))}}
   {\seqq{ce' = ce[(\mbox{``AggrType''}, type_{ce}(e), \mbox{``Node''})/v]}}
    \quad
    \\[1ex]
    \infer[inst of]
    {\varphi_{InferType, \mbox{inst } v_1 \mbox{ of } v_2} ((ce, cm), (ce', cm'))}
   {\seqq{ce' = ce[(prj_{1,2}(type_{ce}(v_2)), \mbox{``InstanceOf''})/v_1] }}
  \quad
    \\[1ex]
    \infer[$st$: skip, when, connect, invalid statement]
    {\varphi_{InferType, st} ((ce, cm), (ce, cm))}
%   {\seqq{rs(v) \neq \bot}} {\seqq{rs' = rs[s(e)/v]}}
   {\seqq{-}}
  \end{gather*}
%  \vspace{-4mm}
  \caption{Semantics of InferType compilation pass}
  \label{tab:infertype-sem}
\end{table}
}
%%%%%%%%%%%%%%%%%%%%%%%%%%%%%%%
%%%%%%%%%%%%%%%%%%%%%%%%%%%%%%%

%%%%%%%%%%%%%%%%%%%%%%%%%%%%%%%
%%%%%%%%%definition of type_ce(e) omitted %%%%%%
%%%%%%%%%%%%%%%%%%%%%%%%%%%%%%%
\hide{
\begin{table}[htbp]
  \footnotesize
  \begin{gather*}
%    \infer[skip]
%    {\seqq{\varphi_{InferType, \mbox{skip}}((ce, cm), (ce, cm))}}
%    {\seqq{-}}
%    \quad
    \infer[UInt const]
%    {\seqq{}}
    {\seqq{type_{ce}(\uint{n}(c)) = (\mbox{``Gtyp''}, \uint{n})}}
   {\seqq{-}}
   \quad
    \infer[SInt const]
    {\seqq{type_{ce}(\sint{n}(c)) = (\mbox{``Gtyp''}, \sint{n})}}
   {\seqq{-}}
    \\[1ex]
    \inferii[mux]
    {type_{ce}(mux(e_0, e_1, e_2)) = maxTypes(t_1, t_2)}
   {\seqq{type_{ce}(e_1) = t_1}} {\seqq{type_{ce}(e_2) = t_2}}
    \quad
%    \\[1ex]
    \infer[validif]
    {\seqq{type_{ce}(validif(e_1, e_2)) = type_{ce}(e_2)}}
   {\seqq{-}}
    \quad
    \\[1ex]
    \inferii[$bop$=add, sub, $\theta=\mbox{\prettyll{UInt}}, \mbox{\prettyll{SInt}}$]
    {type_{ce}(bop(e_1, e_2)) = (\mbox{``Gtype''}, \theta\mbox{<}\max(n_1, n_2)+1\mbox{>})}
   {\seqq{type_{ce}(e_1) = (\mbox{``Gtype''}, \theta\mbox{<}n_1\mbox{>})}} {\seqq{type_{ce}(e_2) =  (\mbox{``Gtype''}, \theta\mbox{<}n_2\mbox{>})}}
    \quad
    \\[1ex]
    \inferii[mul, $\theta=\mbox{\prettyll{UInt}}, \mbox{\prettyll{SInt}}$]
    {type_{ce}(mul(e_1, e_2)) = (\mbox{``Gtype''}, \theta\mbox{<}n_1+ n_2\mbox{>})}
   {\seqq{type_{ce}(e_1) = (\mbox{``Gtype''}, \theta\mbox{<}n_1\mbox{>})}} {\seqq{type_{ce}(e_2) =  (\mbox{``Gtype''}, \theta\mbox{<}n_2\mbox{>})}}
    \quad
    \\[1ex]
    \inferii[div, UInt]
    {type_{ce}(div(e_1, e_2)) = (\mbox{``Gtype''}, \uint{n_1})}
   {\seqq{type_{ce}(e_1) = (\mbox{``Gtype''}, \uint{n_1})}} {\seqq{type_{ce}(e_2) =  (\mbox{``Gtype''}, \uint{n_2})}}
    \quad
    \\[1ex]
    \inferii[div, SInt]
    {type_{ce}(div(e_1, e_2)) = (\mbox{``Gtype''}, \sint{n_1+1})}
   {\seqq{type_{ce}(e_1) = (\mbox{``Gtype''}, \sint{n_1})}} {\seqq{type_{ce}(e_2) =  (\mbox{``Gtype''}, \sint{n_2})}}
    \quad
    \\[1ex]
  \end{gather*}
%  \vspace{-4mm}
  \caption{Definition of $type_{ce}(e)$}
  \label{tab:type-hexp}
\end{table}

$type_{ce}(e)$ is defined as follows.
\begin{itemize}
\item $type_{ce}(\uint{n}(c)) = (\mbox{``Gtyp''}, \mbox{``UInt''})$,
%
\item $type_{ce}(\sint{n}(c)) = (\mbox{``Gtyp''}, \mbox{``SInt''})$,
%
\item let $type_{ce}(e_1) = t_1$ and $type_{ce}(e_2) = t_2$, then $type_{ce}(mux(e_0, e_1, e_2)) = maxTypes(t_2, t_3)$,
%
\item $type_{ce}(validif(e_1, e_2)) = type_{ce}(e_2)$,
%
\item let $type_{ce}(e_1) =t_1$ and $type_{ce}(e_2) = t_2$, then
\begin{itemize}
\item $type_{ce}(add(e_1, e_2)) =\theta \bwidth{\max(n_1, n_2)+1}$ and $type_{ce}(sub(e_1, e_2)) = \theta \bwidth{\max(n_1, n_2)+1}$,
%
\item $type_{ce}(mul(e_1, e_2)) = \theta \bwidth{n_1+n_2}$,
%
\item $type_{ce}(div(e_1, e_2)) = UInt \bwidth{n_1}$ if $\theta = UInt$, and $type_{ce}(div(e_1, e_2)) = SInt \bwidth{n_1+1}$ otherwise,
%
\item $type_{ce}(mod(e_1, e_2)) = \theta \bwidth{\min(n_1,n_2)}$,
%
\item $type_{ce}(cp(e_1, e_2)) = UInt \bwidth{1}$, where $cp \in \{lt, leq, gt, geq, eq, neq\}$,
%
\item $type_{ce}(dshl(e_1, e_2)) = \theta \bwidth{n_1+2^{n_2}-1}$,
%
\item $type_{ce}(dshr(e_1, e_2)) = \theta \bwidth{n_1}$,
%
\item $type_{ce}(and(e_1, e_2)) = type_{ce}(or(e_1, e_2)) =  te(xor(e_1, e_2)) = \theta \bwidth{\max(n_1, n_2)}$,
%
\item $type_{ce}(cat(e_1, e_2)) = \theta \bwidth{n_1 + n_2}$,
\end{itemize}
%
\item let $te(e_1) = \theta \bwidth{n_1}$, where $\theta \in \{UInt, SInt\}$, then
\begin{itemize}
\item $te(andr(e_1)) = te(orr(e_1)) =  te(xorr(e_1)) = \uint{1}$,
%
\item $te(neg(e_1)) = \sint{n_1 +1}$, $te(not(e_1)) =  \uint{n_1}$,
%
\item $te(cvt(e_1)) = \sint{n_1 +1}$ if $\theta = UInt$, and $te(cvt(e_1)) =  \sint{n_1}$ otherwise,
\end{itemize}
%
\item let $te(e_1) = \theta \bwidth{n_1}$, where $\theta \in \{UInt, SInt\}$, then
\begin{itemize}
\item $te(pad(e_1, n)) = \uint{n}$,
%
\item $te(shl(e_1,n)) = \theta \bwidth{n_1 + n}$,
%
\item $te(shr(e_1, n)) = \theta \bwidth{n_1-n}$ if $n_1 > n$, and $te(shr(e_1, n)) =  \theta \bwidth{1}$ otherwise,
%
\item $te(head(e_1, n)) = \uint{n}$, $te(tail(e_1, n)) = \uint{n_1-n}$,
%
\item $te(bits(e_1, n, n')) = \uint{n'-n+1}$.
\end{itemize}
%
\end{itemize}
}

%%%%%%%%%%%%%%%%%%%%%%%%%%%%%%%%%%%%%%%%%%%%%%%%%%%%%%%%%%%%%%%%%%%%%%%%%%
